\exercisesection{Banach 空间自同态的谱}

在下面的习题中，所有 Banach 空间都取在 C 上。

\begin{enumerate}
\def\labelenumi{\arabic{enumi})}
\tightlist
\item
  设 \(\mathrm{E} = \mathcal{C}([0,1])\)，设 \(u\) 为由下式定义的 \(\mathrm{E}\) 的自同态
\end{enumerate}

\[
u (f) (t) = \int_ {0} ^ {t} f (t) d t.
\]

\begin{enumerate}
\def\labelenumi{\alph{enumi})}
\item
  \(u\) 的谱半径为零。（估计 \(u^n\) 的范数。）
\item
  \(u\) 的谱仅由 0 组成，且对每个 \(n \geqslant 1\)，\(u^{n}\) 的核仅由 0 组成。
\end{enumerate}

（特别地，0 是 \(u\) 的谱的一个孤立点，但诸 \(u^n\)（\(n \geqslant 1\)）的核的并集的闭包并不等于谱子空间 \(\mathrm{E}_0(x)\)。）

\begin{enumerate}
\def\labelenumi{\arabic{enumi})}
\setcounter{enumi}{1}
\item
  存在一个 Banach 空间 E 及 E 的一个自同态 u，使得 0 是 u 的特征值，但 0 在 \(\mathrm{Sp}(u)\) 中不是孤立点。
\item
  给出一个 Hilbert 空间 E 与一个自同态 \(u \in \mathcal{L}(\mathrm{E})\) 的例子，使得 0 是 u 的特征值，但不是 \(u^{*}\) 的特征值。
\item
  设 \(E = \ell^{2}(\mathbf{Z})\)，设 \(u\) 为满足 \(u((\lambda_{n})) = (\lambda_{n+1})\) 的自同态。
\end{enumerate}

\begin{enumerate}
\def\labelenumi{\alph{enumi})}
\item
  证明 \(u\) 是酉自同态。
\item
  计算 \(u\) 的谱。
\item
  \(u\) 的有限谱是空集。
\end{enumerate}

\begin{enumerate}
\def\labelenumi{\arabic{enumi})}
\setcounter{enumi}{4}
\tightlist
\item
  设 \(s\) 为 \(\ell^2 (\mathbf{N})\) 的由下式确定的自同态
\end{enumerate}

\[
s (a _ {0}, a _ {1}, \dots , a _ {n}, \dots) = (a _ {1}, \dots , a _ {n}, \dots).
\]

证明 \(s\) 的数值域是 C 中的开单位圆盘 \(\Delta\)，而 \(s\) 的谱是 C 中的闭单位圆盘。求出 \(s\) 的特征值。对 \(s^{*}\)（\(s\) 的伴随自同态）作同样的讨论。

\begin{enumerate}
\def\labelenumi{\arabic{enumi})}
\setcounter{enumi}{5}
\item
  设 \(u \in \mathcal{L}(\mathbf{C}^n)\) 为迹为零的自同态。存在 \(\mathbf{C}^n\) 的一个标准正交基 B，使得 \(u\) 在基 B 下的矩阵的所有对角元均为零。（证明 0 属于 \(u\) 的数值域。）
\item
  设 \(u\) 为复 Hilbert 空间 E 的自同态，使得 \(u\) 的谱包含于 C 中的开单位圆盘。设 F 为由 E 的元素组成的平方可和序列所构成的 Hilbert 空间 \(\ell_{\mathrm{E}}^{2}(\mathbf{N})\)（EVT，V，第 18 页，例），并设 \(s_{\mathrm{F}}\) 为 F 的由下式确定的自同态
\end{enumerate}

\[
s _ {\mathrm{F}} (a _ {0}, a _ {1}, \ldots , a _ {n}, \ldots) = (a _ {1}, \ldots , a _ {n} \ldots).
\]

存在 F 的一个闭子空间 H 与一个连续的双射线性映射 \(v: E \rightarrow H\)，使得 \(u = v^{-1}s_{F}v\)。

\begin{enumerate}
\def\labelenumi{\arabic{enumi})}
\setcounter{enumi}{7}
\tightlist
\item
  设 \(u\) 为复 Hilbert 空间 E 的一个自同态。设 \(S_u\) 为 E 的形如 \(wuw^{-1}\) 的自同态所成的集合，其中 \(w\) 取遍 E 的可逆自同态。
\end{enumerate}

\begin{enumerate}
\def\labelenumi{\alph{enumi})}
\tightlist
\item
  我们有
\end{enumerate}

\[
\varrho (u) = \inf _ {v \in \mathrm{S} _ {u}} \| v \|.
\]

（可设 \(\varrho(u) \leqslant 1\)；对自同态 \(u_{\varepsilon} = (1 + \varepsilon)^{-1} u\)（\(\varepsilon > 0\)）应用前一个习题。）

\begin{enumerate}
\def\labelenumi{\alph{enumi})}
\setcounter{enumi}{1}
\tightlist
\item
  \(u\) 的谱的凸包等于诸子集 \(\iota(v)\)（\(v \in S_u\)）的交（"Hildebrandt 定理"）。
\end{enumerate}

¶ 9) 设 E 为复 Banach 空间。

\begin{enumerate}
\def\labelenumi{\alph{enumi})}
\item
  设 \(u \in \mathcal{L}(\mathrm{E})\)。考察 \(\gamma_{u}\) 与 \(\delta_{u}\)，即 \(\mathcal{L}(\mathrm{E})\) 的由 \(\gamma_{u}(v) = u \circ v\) 与 \(\delta_{u}(v) = v \circ u\) 定义的自同态。我们有 \(\operatorname{Sp}(u) = \operatorname{Sp}(\gamma_{u}) = \operatorname{Sp}(\delta_{u})\)，其中后两个谱是相对于代数 \(\mathcal{L}(\mathcal{L}(\mathrm{E}))\) 取的。
\item
  对每个函数 \(f\)，只要它在 \(u\) 的谱的某邻域内全纯，就有 \(\boldsymbol{\gamma}_{f(u)} = f(\boldsymbol{\gamma}_{u})\) 与 \(\boldsymbol{\delta}_{f(u)} = f(\boldsymbol{\delta}_{u})\)。
\item
  设 \(u_{1}\) 与 \(u_{2}\) 属于 \(\mathcal{L}(\mathrm{E})\)。设 \(\varrho\) 为映射 \(v \mapsto u_{1} \circ v \circ u_{2}\)，它是 \(\mathcal{L}(\mathrm{E})\) 的一个自同态。它的谱等于 \(\operatorname{Sp}(u_{1})\operatorname{Sp}(u_{2}) \subset \mathbf{C}\)。
\end{enumerate}

\(d)\) 设 \(u_{1}\) 与 \(u_{2}\) 属于 \(\mathcal{L}(\mathrm{E})\)。设 \(\varrho\) 为映射 \(v \mapsto u_{1} \circ v \circ u_{2}\)，它是从 E 到 E 的 Hilbert--Schmidt 映射所成的空间 \(\mathcal{L}_{2}(\mathrm{E})\) 的一个自同态。它的谱等于 \(\mathrm{Sp}(u_{1})\mathrm{Sp}(u_{2}) \subset \mathbf{C}\)。（应用类似的方法。）

\begin{enumerate}
\def\labelenumi{\arabic{enumi})}
\setcounter{enumi}{9}
\item
  设 E 为复 Hilbert 空间。设 \(x_1\) 与 \(x_2\) 为 E 的自同态。证明 \(x_1 \otimes x_2\) 作为 E \(\widehat{\otimes}_2\) E 的自同态，其谱等于 \(\mathrm{Sp}(x_1) \mathrm{Sp}(x_2)\)。（把 \(\overline{\mathrm{E}} \widehat{\otimes}_2\) E 与从 E 到 E 的 Hilbert--Schmidt 算子所成的空间等同起来，参见 EVT 5，第 52 页，定理 1，b)，并应用前一个习题。）
\item
  设 \(u\) 为复 Hilbert 空间 \(E\) 的自同态，并设 \((j, |u|)\) 为它的极分解。
\end{enumerate}

\begin{enumerate}
\def\labelenumi{\alph{enumi})}
\item
  为使 \(j\) 与 \(|u|\) 交换，充要条件是 \(u\) 与 \(u^{*}u\) 交换。
\item
  证明此时 \(\|u^{*}(x)\|\leqslant\|u(x)\|\) 对一切 \(x\in E\) 成立，并且 \(\varrho(u)=\|u\|\)。
\item
  设 \(E = \ell^{2}(\mathbf{N})\)，并设 \(u\) 为由下式给出的 E 的自同态
\end{enumerate}

\[
u ((a _ {0}, a _ {1}, \dots)) = (0, a _ {0}, a _ {1}, \dots).
\]

它不是正规的。计算 u 的极分解 \((j, |u|)\)；\(j\) 与 \(|u|\) 交换。

\begin{enumerate}
\def\labelenumi{\arabic{enumi})}
\setcounter{enumi}{11}
\tightlist
\item
  设 \(u\) 为 Banach 空间 E 的一个自同态。用 \(\sigma_{a}(u)\) 表示满足下列条件的 \(\lambda\in\mathbf{C}\) 所成的集合：对每个 \(\alpha\in\mathbf{R}_{+}^{*}\)，存在 \(x\in\mathrm{E}\) 使 \(\|\lambda x-u(x)\|<\alpha\|x\|\)。
\end{enumerate}

\begin{enumerate}
\def\labelenumi{\alph{enumi})}
\item
  证明 \(\sigma_{a}(u)\) 是 \(\sigma(u)\) 的一个闭子集。
\item
  证明 \(\sigma_{a}(u)\) 包含 \(\sigma(u)\) 的边界。
\end{enumerate}

\begin{enumerate}
\def\labelenumi{\arabic{enumi})}
\setcounter{enumi}{12}
\tightlist
\item
  设 E 为复 Hilbert 空间。设 N 为从 \(\mathcal{L}(\mathrm{E})\) 到 \(\mathfrak{P}(\mathbf{C})\) 中的一个映射，使得
\end{enumerate}

\begin{enumerate}
\def\labelenumi{\alph{enumi})}
\item
  \(N(u)\) 对每个 \(u \in \mathcal{L}(\mathrm{E})\) 都是紧的；
\item
  \(\mathrm{N}(\alpha u + \beta 1_{\mathrm{E}}) = \alpha \mathrm{N}(u) + \beta\) 对一切 \(u \in \mathcal{L}(\mathrm{E})\) 与每个 \((\alpha, \beta) \in \mathbf{C}^{2}\) 成立；
\item
  \(\mathrm{N}(u)\subset\{z\in\mathbf{C}\mid\mathcal{R}(z)\geqslant0\}\) 当且仅当 \(u+u^{*}\) 是正的。
\end{enumerate}

  那么 \(\mathrm{N}(u)=\overline{\iota(u)}\) 对一切 \(u\in\mathcal{L}(\mathrm{E})\) 成立。（证明 \(u\mapsto\overline{\iota(u)}\) 具有所指出的性质；设 \(N_{1}\) 与 \(N_{2}\) 为满足这些性质的映射；假设 \(\mathrm{N}_{1}(u)\) 不包含于 \(\mathrm{N}_{2}(u)\)，求 C 中的 \(\alpha\) 与 \(\beta\)，使 \(\mathrm{N}_{2}(\alpha u+\beta)\) 包含于半平面 \(\mathcal{R}(z)\geqslant0\) 而 \(\mathrm{N}_{1}(\alpha u+\beta)\) 不然，再利用 c) 得出矛盾。）

\begin{enumerate}
\def\labelenumi{\arabic{enumi})}
\setcounter{enumi}{13}
\item
  给出一个复 Hilbert 空间及 \(u \in \mathcal{L}(\mathrm{E})\) 的例子，使 \(\iota(u)\) 不是闭的。
\item
  C 的非空紧子集所成的集合 X 赋以 TG，IX，第 91 页，习题 6 中定义的距离。设 E 为复 Hilbert 空间。
\end{enumerate}

\begin{enumerate}
\def\labelenumi{\alph{enumi})}
\item
  映射 \(u \mapsto \overline{\iota(u)}\) 从 \(\mathcal{L}(\mathrm{E})\) 到 X 是连续的。若 E 是无穷维的，则当 \(\mathcal{L}(\mathrm{E})\) 赋以单收敛拓扑时，它不连续。
\item
  映射 \(u \mapsto \operatorname{Sp}(u)\) 从 \(\mathcal{L}(\mathrm{E})\) 到 X 不连续。
\end{enumerate}

\begin{enumerate}
\def\labelenumi{\arabic{enumi})}
\setcounter{enumi}{15}
\tightlist
\item
  设 X 为局部紧空间，\(\mu\) 为 X 上总质量为 1 的正测度，\(f: X \to X\) 为满足 \(f(\mu) = \mu\) 的映射。
\end{enumerate}

\begin{enumerate}
\def\labelenumi{\alph{enumi})}
\item
  映射 \(\varphi \mapsto \varphi \circ f\) 诱导出酉自同态 \(u_{f}\)，它定义在 \(\mathrm{L}^2 (\mathrm{X},\mu)\) 上。
\item
  实数 1 是 \(u_f\) 的重数为 1 的特征值，其充要条件是：对 X 的每个满足 \(f(A) = A\) 的可测子集 A，有 \(\mu(A) \in \{0,1\}\)。这时称 \(f\) 是 \(\mu\)-遍历的。
\item
  设 \(\mathrm{E}\subset\mathrm{L}^{2}(\mathrm{X},\mu)\) 为 \(u_{f}\) 关于特征值 1 的特征子空间。序列 \((v_{\mathrm{N}})_{\mathrm{N}\in\mathbf{N}}\) 定义为
\end{enumerate}

\[
v _ {\mathrm{N}} = \frac {1}{\mathrm{N}} \sum_ {n = 0} ^ {\mathrm{N-1}} u _ {f} ^ {n}
\]

在赋以单收敛拓扑的空间 \(\mathcal{L}(\mathrm{L}^{2}(\mathrm{X},\mu))\) 中收敛到 E 上的正交投影（"von Neumann 遍历定理"）。（注意 E 的正交补是 \(u_{f}-1_{\mathrm{E}}\) 的像的闭包；然后验证 \(v_{\mathrm{N}}(\varphi)\to0\) 当 \(\varphi\in\mathrm{E}^{\circ}\) 时成立。）

\(d)\) 定义 \(\widetilde{u}_f\)，即 \(\mathrm{L}^1 (\mathrm{X},\mu)\) 的一个如 \(a)\) 中那样的自同态，并令

\[
\widetilde {v} _ {\mathrm{N}} = \frac {1}{\mathrm{N}} \sum_ {n = 0} ^ {\mathrm{N-1}} \widetilde {u} _ {f} ^ {n}.
\]

对每个 \(\varphi\in\mathrm{L}^{1}(\mathrm{X},\mu)\)，序列 \((\widetilde{v}_{\mathrm{N}}(\varphi))_{\mathrm{N}\in\mathbf{N}}\) 在 \(\mathrm{L}^{1}(\mathrm{X},\mu)\) 中收敛到元素 \(\widetilde{\varphi}\)，它属于 \(\widetilde{u}_{f}\) 关于特征值 1 的特征子空间。（先考虑 \(\varphi\) 有界的情形，并应用前一个问题。）

\begin{enumerate}
\def\labelenumi{\alph{enumi})}
\setcounter{enumi}{4}
\tightlist
\item
  设 I 为有限集，X = I \(^{Z}\) 赋以 I 上的测度 \(\mu_n\) 的乘积测度，其中 \(\mu_n(i) = 1/\text{Card(I)}\) 对一切 \(i \in \text{I}\) 成立。令 \(f((i_n)) = (i_{n+1})\) 对一切 \((i_n) \in \text{X}\) 成立。那么 \(f(\mu) = \mu\)，且 \(f\) 是 \(\mu\)-遍历的（“Bernoulli 移位”）。
\end{enumerate}

\(f)\) 设 \(\mathrm{X} = \mathbf{R} / \mathbf{Z}\)，\(\mu\) 为 \(\mathrm{X}\) 上的规范化 Haar 测度。设 \(a \in \mathrm{X}\)，令 \(f(x) = x + a\)。那么 \(f(\mu) = \mu\)，且 \(f\) 是 \(\mu\)-遍历的当且仅当 \(a\) 是某个无理数的类。

\exercisesection{紧空间上连续函数的代数}

下面所有的 Banach 空间与所有的 Banach 代数都是复的。

\begin{enumerate}
\def\labelenumi{\arabic{enumi})}
\tightlist
\item
  设 X 为拓扑空间，\(\mathcal{C}_{b}(\mathrm{X})\) 为 X 上有界连续复函数所成的星代数，并设 \(\widehat{\mathrm{X}} = \mathrm{X}(\mathcal{C}_{b}(\mathrm{X}))\)。
\end{enumerate}

\begin{enumerate}
\def\labelenumi{\alph{enumi})}
\tightlist
\item
  设 A 为 \(\mathcal{C}_{b}(\mathrm{X})\) 的一个幺 Banach 子代数，设 \(R_{A}\) 为 X 上的如下等价关系：x 与 \(x'\) 等价，如果 \(f(x) = f(x')\) 对每个函数 \(f \in A\) 成立。X 的紧子集 K 的 \(R_{A}\)-饱和集 S 是闭的。（设 \(x \in \overline{S}\)；对一切 \(f \in A\) 及每个 \(\varepsilon > 0\)，设 \(V_{f,\varepsilon}\) 为由那些 \(x' \in X\)（满足 \(|f(x) - f(x')| \leqslant \varepsilon\)）所组成的集合；若 \(f_{1}, \ldots, f_{n} \in A\) 且 \(\varepsilon_{1}, \ldots, \varepsilon_{n} > 0\)，则有
\end{enumerate}

\[
\mathrm{V} _ {f _ {1}, \varepsilon_ {1}} \cap \dots \cap \mathrm{V} _ {f _ {n}, \varepsilon_ {n}} \cap \mathrm{K} \neq \varnothing ,
\]

因此

\[
\bigcap_{\substack{f\in \mathrm{A}\\ \varepsilon >0}}\mathrm{V}_{f,\varepsilon}\cap \mathrm{K}\neq \varnothing
\]

从而 \(x \in S\)。）

\begin{enumerate}
\def\labelenumi{\alph{enumi})}
\setcounter{enumi}{1}
\item
  假设 X 是正规的。设 R 为 X 上的闭等价关系，并设 \(A = A_{R}\) 为由在模 R 的诸类上取常值的函数所组成的 \(\mathcal{C}_{b}(X)\) 的子代数。证明 \(R = R_{A}\)（利用 TG，IX，第 105 页，习题 19，a）。
\item
  假设 X 是紧的。\(R \mapsto A_{R}\) 是从 X 上分离等价关系所成的集合到 \(\mathcal{C}(X)\) 的星子代数所成的集合上的一个双射（设 A 为 \(\mathcal{C}(X)\) 的一个星子代数。证明 \(R_{A}\) 是分离的，并对由 A 过渡到商而得到的 \(X/R_{A}\) 上的连续函数代数应用 Weierstrass–Stone 定理。）
\end{enumerate}

\begin{enumerate}
\def\labelenumi{\arabic{enumi})}
\setcounter{enumi}{1}
\tightlist
\item
  设 X 为紧拓扑空间，R 为 X 上的分离等价关系，B 为 \(\mathcal{C}(\mathrm{X})\) 的包含 \(\mathrm{A}_{\mathrm{R}}\) 的闭子代数（习题 1 的记号）。若 \(f \in \mathcal{C}(\mathrm{X})\) 在每个 R-类 \(c\) 上与 B 的一个元素 \(g_{c}\) 相一致，则 \(f \in \mathrm{B}\)。（设 \(\varepsilon > 0\)。存在 \(\mathrm{V}_{c}\)，即 \(c\) 的一个饱和开邻域，使 \(|f - g_{c}| \leqslant \varepsilon\) 在 \(\mathrm{V}_{c}\) 上成立。设 \(\mathrm{V}_{c_{1}}, \ldots, \mathrm{V}_{c_{n}}\) 为 X 的一个覆盖。设 \((u_{1}, \ldots, u_{n})\) 为从属于 \((\mathrm{V}_{c_{1}}, \ldots, \mathrm{V}_{c_{n}})\) 且由 \(\mathrm{A}_{\mathrm{R}}\) 中的函数组成的单位分解。那么
\end{enumerate}

\[
g = u _ {1} g _ {c _ {1}} + \dots + u _ {n} g _ {c _ {n}} \in \mathrm{B},
\]

且 \(|f - g| \leqslant \varepsilon\)。）

\begin{enumerate}
\def\labelenumi{\arabic{enumi})}
\setcounter{enumi}{2}
\tightlist
\item
  设 X 为紧拓扑空间，且对每个 \(x \in X\)，设 \(A_{x}\) 为交换幺 Banach 代数，用 \(e_{x}\) 表示其单位元。我们假设 \(\|e_{x}\| = 1\)。设 B 为诸 \(A_{x}\)（\(x \in X\)）的乘积 Banach 代数。设 C 为 B 的具有下列性质的 Banach 子代数：
\end{enumerate}

\begin{enumerate}
\def\labelenumi{(\roman{enumi})}
\item
  \(e = (e_x)_{x \in \mathbf{X}} \in \mathbf{C};\)
\item
  若 \(a = (a_x) \in \mathrm{C}\) 且 \(f \in \mathcal{C}(\mathrm{X})\)，则函数 \(fa: x \mapsto f(x)a_x\) 是 \(\mathrm{C}\) 的一个元素；
\item
  对每个 \(a = (a_x) \in \mathrm{C}\)，函数 \(x \mapsto \|a_x\|\) 是上半连续的。
\end{enumerate}

设 \(\widehat{\mathbf{X}}\) 为诸 \(\mathsf{X}(\mathsf{A}_x)\) 的不交并。对每个 \(x_0 \in \mathbf{X}\) 与每个 \(\chi \in \mathsf{X}(\mathsf{A}_{x_0})\)，设 \(\varphi(\chi)\) 为由 \(\varphi(\chi)((a_x)_{x \in \mathbf{X}}) = \chi(a_{x_0})\) 定义的 C 的特征标。 证明 \(\varphi\) 给出从 \(\widehat{\mathbf{X}}\) 到 \(\mathsf{X}(\mathbf{C})\) 上的一个双射。（设 \(\xi \in \mathsf{X}(\mathbf{C})\)。由于 \(\mathcal{C}(\mathbf{X})\) 可等同于 C 的一个闭子代数，\(\xi\) 定义了 \(\mathcal{C}(\mathbf{X})\) 的一个特征标，于是存在 \(x_0 \in \mathbf{X}\)。设 \(a = (a_x) \in \mathbf{C}\) 与 \(b = (b_x) \in \mathbf{C}\)，满足 \(a_{x_0} = b_{x_0}\)；设 \(\varepsilon > 0\)；\(\|a_x - b_x\| < \varepsilon\) 在 \(x_0\) 的某邻域 U 上成立；设 \(f \in \mathcal{C}(\mathbf{X})\) 满足 \(0 \leqslant f \leqslant 1\)，在 \(x_0\) 处等于 1 且在 U 外等于 0；我们有 \(\|fa - fb\| \leqslant \varepsilon\)，从而 \(|\xi(fa - fb)| \leqslant \varepsilon\)，从而 \(|\xi(a) - \xi(b)| \leqslant \varepsilon\)，从而 \(\xi(a) = \xi(b)\)。由此推出存在 \(\chi \in \mathsf{X}(\mathsf{A}_{x_0})\) 使 \(\xi = \varphi(\chi)\)。

\begin{enumerate}
\def\labelenumi{\arabic{enumi})}
\setcounter{enumi}{3}
\tightlist
\item
  设 X 为紧拓扑空间，使得 Banach 代数 \(\mathcal{C}(\mathrm{X})\) 由一个幂等元序列 \((j_n)\) 生成。函数
\end{enumerate}

\[
\omega \mapsto \sum_ {n = 1} ^ {\infty} \frac {1}{3 ^ {n}} (2 j _ {n} (\omega) - 1)
\]

分离 X 的点。它生成 Banach 代数 \(\mathcal{C}(\mathrm{X})\)。

\begin{enumerate}
\def\labelenumi{\arabic{enumi})}
\setcounter{enumi}{4}
\tightlist
\item
  设 A 为从 {[}0,1{]} 到 C 的、在 {[}0,1{]} 上具有直到 n 阶连续导数的函数所成的 Banach 代数（I，第 18 页，例 4）。设 \(\omega \in [0,1] = \mathsf{X}(\mathsf{A})\)。
\end{enumerate}

\begin{enumerate}
\def\labelenumi{\alph{enumi})}
\item
  A 中使 \(V(I) = \{\omega\}\) 成立的最小闭理想 I 是由这样的那些 \(f \in A\) 组成的集合：\(f, f', \ldots, f^{(n)}\) 在 \(\omega\) 处皆为 0。（利用 I，第 94 页，命题 5。）
\item
  若 \(g \in A\)，其典范像在 A/I 中的范数等于
\end{enumerate}

\[
\sum_ {k = 0} ^ {n} \frac {| g ^ {(k)} (\omega) |}{k !}.
\]

\begin{enumerate}
\def\labelenumi{\alph{enumi})}
\setcounter{enumi}{2}
\tightlist
\item
  代数 A/I 同构于 \(\mathbf{C}[\mathrm{X}]/(\mathrm{X}^{n+1})\)；A 中包含 I 的唯一极大理想是诸 \(f \in A\) 中使 \(f(\omega) = 0\) 者所成的集合。
\end{enumerate}

\begin{enumerate}
\def\labelenumi{\arabic{enumi})}
\setcounter{enumi}{5}
\tightlist
\item
  设 \(\Delta\) 为由 \(|\zeta| \leqslant 1\) 给出的 \(\mathbf{C}\) 中的圆盘，A 为在 \(\Delta\) 上连续且在 \(\Delta\) 的内部解析的复值函数所成的 Banach 代数（I，第 20 页，例 9）。
\end{enumerate}

\begin{enumerate}
\def\labelenumi{\alph{enumi})}
\tightlist
\item
  若 \(x \in A\)，则 \(x\) 在 A 中是下列函数的极限：
\end{enumerate}

\[
\zeta \mapsto x _ {n} (\zeta) = x \left(\frac {n \zeta}{n + 1}\right).
\]

\(\Delta\) 的恒等映射是 A 的一个生成元。

\begin{enumerate}
\def\labelenumi{\alph{enumi})}
\setcounter{enumi}{1}
\item
  空间 X(A) 典范地等同于 \(\Delta\)。（利用 I，第 142 页，命题 1 的断言 g）。）
\item
  设 S 为 \(\Delta\) 的一个闭子集，并把它等同于 \(\mathsf{X}(\mathsf{A})\) 的一个闭子集。若 S 有一个非孤立点 \(\zeta\)，满足 \(|\zeta| < 1\)，则 S 对 Jacobson 拓扑的闭包是 \(\Delta\)。
\end{enumerate}

\(d)\) 对 \(x \in \mathrm{A}\)，设 \(x^{*}(\zeta) = x(\overline{\zeta})\)。则 \(x \mapsto x^{*}\) 是 A 的一个等距对合。A 的 Hermite 特征标由 \(\Delta\) 的实元素给出。

\begin{enumerate}
\def\labelenumi{\alph{enumi})}
\setcounter{enumi}{4}
\tightlist
\item
  映射 \(x \mapsto x|U\) 是从 A 到 \(\mathrm{P}(\mathbf{U})\) 上的一个同构。\(\mathcal{C}_{\mathbf{R}}(\mathbf{U})\) 中的每个函数都是 \(\mathrm{P}(\mathbf{U})\) 中函数的实部的一致极限。
\end{enumerate}

\begin{enumerate}
\def\labelenumi{\arabic{enumi})}
\setcounter{enumi}{6}
\tightlist
\item
  设 \(A_{1}\)（相应地，\(A_{2}\)）为习题 5（相应地，习题 6）中所考虑的幺 Banach 代数。设 \(A = A_{1} \times A_{2}\)。证明 A 是半单的，但 A 的 Gelfand 变换的像在 \(\mathcal{C}(\mathsf{X}(\mathsf{A}))\) 中既非闭也非稠密。
\end{enumerate}

¶ 8) 设 X 为紧拓扑空间，B 为 \(\mathcal{C}(\mathrm{X})\) 的一个分离 X 的点的幺 Banach 子代数（带诱导范数）。把 X 等同于 \(\mathrm{X(B)} = \mathrm{X}'\) 的一个闭子集。设 \(\mathrm{B}^*\) 为 B 中可逆元所成的集合。若诸函数 \(\log(|f|)\)（\(f \in \mathrm{B}^*\)）所成的集合在 \(\mathcal{C}_{\mathbf{R}}(\mathrm{X})\) 中稠密，则称 B 为对数模代数（logmodular）。（特别地，当诸函数 \(\mathcal{R}(f)\)（\(f \in \mathrm{B}\)）所成的集合在 \(\mathcal{C}_{\mathbf{R}}(\mathrm{X})\) 中稠密时即是这种情形，因为 \(\log(|e^f|) = \mathcal{R}(f)\)。）以下假设 B 是对数模代数。

\begin{enumerate}
\def\labelenumi{\alph{enumi})}
\tightlist
\item
  对每个 \(\chi \in \mathrm{X}'\)，在 X 上存在唯一的测度 \(\mu_{\chi} \geqslant 0\) 使得
\end{enumerate}

\[
\log (| \chi (f) |) = \int_ {\mathrm{X}} \log (| f (\omega) |) d \mu_ {\chi} (\omega)
\]

对所有 \(f \in \mathrm{B}^*\) 成立，从而 \(\chi(f) = \int_{\mathrm{X}} f(\omega) d\mu_{\chi}(\omega)\) 对所有 \(f \in \mathrm{B}\) 成立。\(b)\) 我们有

\[
\log | \chi (f) | \leqslant \int_ {\mathrm{X}} \log | f (\omega) | d \mu_ {\chi} (\omega)
\]

对所有 \(f \in \mathrm{B}\) 成立，并且条件 \(\chi(f) = \int_{\mathrm{X}} f(\omega) d\mu_{\chi}(\omega)\)（对所有 \(f \in \mathrm{B}\)）唯一地确定 \(\mu_{\chi}\)。

\begin{enumerate}
\def\labelenumi{\alph{enumi})}
\setcounter{enumi}{2}
\tightlist
\item
  设 \(\chi \in X'\)，并设 \(\mu\) 为 \(X\) 上的一个正测度。把 μ 写成 \(\mu = \mu_1 + \mu_2\)，其中 \(\mu_1\) 关于 \(\mu_{\chi}\) 绝对连续，而 \(\mu_2\) 与 \(\mu_{\chi}\) 互相奇异。设 I 为 \(\chi\) 的核。则
\end{enumerate}

\[
\inf _ {f \in \mathrm{I}} \int_ {\mathrm{X}} | 1 - f | ^ {2} d \mu = \inf _ {f \in \mathrm{I}} \int_ {\mathrm{X}} | 1 - f | ^ {2} d \mu_ {1}.
\]

\(d)\) 设 \(\chi\in\mathrm{X}^{\prime}\)，\(h\) 为 \(\mathcal{L}^{1}(\mu_{\chi})\) 中的一个非负函数。设 I 为 \(\chi\) 的核。则

\[
\inf _ {f \in \mathrm{I}} \int_ {\mathrm{X}} | 1 - f | ^ {2} h d \mu_ {\chi} = \exp \Bigl (\int_ {\mathrm{X}} \log (h) d \mu_ {\chi} \Bigr).
\]

\begin{enumerate}
\def\labelenumi{\alph{enumi})}
\setcounter{enumi}{4}
\tightlist
\item
  对 X 上每个正测度 \(\mu\) 以及每个 \(p \in [1, +\infty[\)，设 \(\mathrm{H}^p(\mu)\) 为 B 在 \(\mathrm{L}^p(\mu)\) 中的典范像的闭包。设 \(\chi \in \mathrm{X}'\)，I 为 \(\chi\) 的核。则 \(\mathrm{L}^2(\mu_\chi)\) 是 \(\mathrm{H}^2(\mu_\chi)\) 与在 \(\mathrm{L}^2(\mu_\chi)\) 中所取的 \(\bar{\mathrm{I}}\)（属于 I 的函数的共轭所成的集合）的典范像的闭包的 Hilbert 和。
\end{enumerate}

\(f)\) 空间 \(\mathrm{H}^1 (\mu_\chi)\) 是诸类 \(\widetilde{h}\)（\(h\in \mathcal{L}^1 (\mu_\chi)\)）所成的集合，这些类满足 \(\int_{\mathrm{X}}fh  d\mu_{\chi} = 0\) 对所有 \(f\in \mathrm{I}\) 成立。

\(g)\) 设 \(h\) 为 \(\mathcal{L}^1 (\mu_\chi)\) 中的一个非负函数。则 \(\widetilde{h} = |\widetilde{f}|\)（其中 \(\widetilde{f}\in \mathrm{H}^1 (\mu_\chi)\) 且 \(\int_{\mathrm{X}}f  d\mu_{\chi}\neq 0\)）成立的充分必要条件是 \(\log (h)\in \mathcal{L}^1 (\mu_\chi)\)。

¶ 9) 设 X 为紧拓扑空间，B 为 \(\mathcal{C}(\mathrm{X})\) 的一个幺 Banach 子代数（带诱导范数）。X 的每个元素 \(x\) 都定义 B 的一个特征标 \(\chi_x\)。若 \(x\) 与 \(x'\) 属于 X，则我们记 \(x \sim x'\)，如果 \(\| \chi_x - \chi_{x'} \| \neq 2\)。a) 设 \(x\) 与 \(x'\) 属于 X。若存在 B 中元素的一个序列 \((f_n)\)，使得 \(\| f_n \| \leqslant 1\)、\(|f_n(x)|\) 趋于 1，并且 \(\liminf_{n \to \infty} |f_n(x) - f_n(x')| > 0\)，则 \(x\not\sim x'\)。（可以假设 \(f_n(x)\) 趋于 1。考虑 \(g_n = (f_n - \lambda_n)(1 - \lambda_n f_n)^{-1}\)，其中 \((\lambda_n)\) 是 \([0,1[\) 中趋于 1 的一个数列。我们有 \(\| g_n \| \leqslant 1\)。若序列 \((\lambda_n)\) 选得适当，则 \(g_n(x)\) 趋于 1 且 \(g_n(x')\) 趋于 -1。）

\begin{enumerate}
\def\labelenumi{\alph{enumi})}
\setcounter{enumi}{1}
\item
  设 R 为 B 中元素的实部所成的集合。若 \(x \sim x'\)，则存在 \(c \in ]0,1[\)，使得 \(f(x) \geqslant cf(x')\) 与 \(f(x') \geqslant cf(x)\) 对 R 中每个 \(f \geqslant 0\) 成立。（只需得到第一个条件；若它不成立，则存在 R 中的序列 \((f_n)\)，满足 \(f_n \geqslant 0\)、\(f_n(x') = 1\)，且 \(f_n(x)\) 趋于 0；设 \(g_n \in B\) 满足 \(\mathcal{R}(g_n) = f_n\)。则 \(e^{-g_n} \in \mathrm{B}\)，\(\| e^{-g_n} \| \leqslant 1\)，\(|e^{-g_n(x')}| = 1 / e\)，而 \(|e^{-g_n(x)}|\) 趋于 1，这与 a) 矛盾。）
\item
  在 X 上存在正测度 \(\alpha\) 与 \(\beta\)，使得 \(f(x)-cf(x')=\alpha(f)\) 与 \(f(x')-cf(x)=\beta(f)\) 对 \(f\in R\) 成立。令 \(\mu_{x}=(1-c^{2})^{-1}(c\beta+\alpha)\)，\(\mu_{x'}=(1-c^{2})^{-1}(c\alpha+\beta)\)，则 \(\mu_{x}(f)=f(x)\) 与 \(\mu_{x'}(f)=f(x')\) 对所有 \(f\in B\) 成立，且 \(c\mu_{x}\leqslant\mu_{x'}\)，\(c\mu_{x'}\leqslant\mu_{x}\)。
\item
  仍设 \(x \sim x'\)。证明：若 \(\mu\) 是 X 上的一个正测度，使得 \(\mu(f) = f(x)\) 对所有 \(f \in B\) 成立，则存在 X 上的一个正测度 \(\mu'\)，使得 \(\mu'(f) = f(x')\) 对所有 \(f \in B\) 成立，且 \(c\mu \leqslant \mu'\) 对某个 \(c \in]0,1[\) 成立。（沿用 c) 的记号，令 \(\mu' = \mu_{x'} - c\mu_x + c\mu\)。）
\item
  关系 \(x \sim x'\) 是 X 上的一个等价关系。（利用 d)。）
\end{enumerate}

¶ 10) 设 X 为紧拓扑空间，A 为 \(\mathcal{C}(\mathrm{X})\) 的一个带诱导范数的幺 Banach 子代数。

\begin{enumerate}
\def\labelenumi{\alph{enumi})}
\item
  X 的子集 K 称为（关于 A 的）反对称子集，如果每个在 K 上取实值的函数 \(f \in A\) 都在 K 上取常值。每个反对称子集都包含于某个极大反对称子集。极大反对称子集是闭的。极大反对称子集所成的集合 \(\mathfrak{R}\) 是 X 的一个划分。
\item
  设 \(A^{\perp}\) 为 A 在 X 上复测度所成的 Banach 空间中的正交子空间。设 \(\mu\) 为 \(A^{\perp}\) 的单位球的一个极点。证明 \(\mu\) 的支集是反对称的。
\item
  设 \(f \in \mathcal{C}(\mathrm{X})\)。若 \(f|K \in A|K\) 对所有 \(K \in \mathfrak{R}\) 成立，则 \(f \in A\)。（应用 b) 与 Krein–Milman 定理。）由此重新导出 Weierstrass–Stone 定理（TG，X，第 36 页，定理 3）。
\end{enumerate}

\(d)\) X 的非空子集 E 称为（关于 A 的）峰集（peak），如果存在 \(f\in \mathrm{A}\) 使得 \(\| f\| = 1\) 且 E 是诸 \(x\in \mathrm{X}\) 中使 \(f(x) = 1\) 者所成的集合。于是可以假设 \(|f(x)| < 1\) 对 \(x\notin \mathrm{E}\) 成立（把 \(f\) 换成 \(\frac{1}{2} (1 + f)\)）。峰集的任何非空可数交仍是峰集。若 \(\mathrm{E}_1,\mathrm{E}_2\) 是峰集，则 \(\mathrm{E}_1\cup \mathrm{E}_2\) 是峰集。（设 \(f_{1},f_{2}\in \mathrm{A}\)，使 \(f_{i} = 1\) 在 \(\mathrm{E}_i\) 上成立，\(|f_i| < 1\) 在 \(\mathrm{X - E}_i\) 上成立；设 \(g_{i} = \frac{1}{4} (1 - f_{i})^{1 / 3}\in \mathrm{A}\)；则 \(1 - g_{1}g_{2} = 1\) 在 \(\mathrm{E}_1\cup \mathrm{E}_2\) 上成立，\(|1 - g_1g_2| < 1\) 在 \(\mathrm{X - (E_1\cup E_2)}\) 上成立。）

\begin{enumerate}
\def\labelenumi{\alph{enumi})}
\setcounter{enumi}{4}
\item
  设 E 是峰集的一个交。设 I 为在 E 上为零的 \(f \in A\) 所成的集合。则 \(f \mapsto f|E\) 定义从 A/I 到 A 中函数在 E 上的限制所成的代数 A\textbar E 上的一个等距。
\item
  设 E 为一个峰集。设 \(E' \subset E\) 为关于 A\textbar E 的一个峰集。则 \(E'\) 是关于 A 的一个峰集。
\item
  X 的每个反对称子集都是峰集的一个交。（利用 d) 与 f)。）
\item
  对 X 的每个反对称子集 K，代数 A\textbar K 是 \(\mathcal{C}(\mathrm{K})\) 的一个闭子代数。（利用 \(e\) 与 \(g\)。）
\item
  设 R 为 A 中元素的实部所成的集合。假设 \(X \in \mathfrak{R}\) 且 R 对乘法稳定。则 X 由单独一点组成。（设 \(x_{0} \in X\)。对 \(u \in R\)，存在唯一的 \(f \in A\) 使得 \(u = \mathcal{R}(f)\) 且 \(f(x_{0}) \in \mathbf{R}\)；令 \(\mathrm{N}(u) = \|f\|\)。则 R 对范数 N 是一个实 Banach 空间；闭图像定理证明 R 中的乘法分别连续，从而连续。由此推出 \(S = R + iR\) 上有一个 Banach 代数结构。设 \(p \in R\) 满足 \(p(x) > 0\) 对所有 \(x \in X\) 成立。证明对 S 的每个特征标 \(\chi\) 有 \(\mathcal{R}(\chi(p)) > 0\)，并由此推出 \(\log(p) \in R\)。假设 \(X \neq \{x_{0}\}\)，可以构造一个 \(f \in A\)，使得在 X 上 \(0 < \mathcal{R}(f) \leqslant 1\)，\(f(x_{0}) \in \mathbf{R}\)，且 \(\|f\|\) 可以任意大。根据前文，存在 \(V \in A\) 使 \(|V|^{2} = \mathcal{R}(f)\)。我们有 \(\|V\| \leqslant 1\)，从而 \(\mathrm{N}(\mathcal{R}(V)) \leqslant 2\)，但
\end{enumerate}

\[
\mathrm{N} (\mathcal {R} (\mathrm{V}) ^ {2}) \geqslant \frac {1}{2} \left(\| \mathrm{V} ^ {2} + f \| - | \mathcal {I} (\mathrm{V}) (x _ {0}) ^ {2} |\right)
\]

可以任意大。）

\begin{enumerate}
\def\labelenumi{\alph{enumi})}
\setcounter{enumi}{9}
\tightlist
\item
  若 R 对乘法稳定，则 A = C(X)。（利用 a)、c)、h)、i)。）因此，若 A \(\neq\) C(X)，则存在 \(u \in R\) 使得 \(u^{2} \notin R\)。
\end{enumerate}

\begin{enumerate}
\def\labelenumi{\arabic{enumi})}
\setcounter{enumi}{10}
\tightlist
\item
  设 X 为一度量空间，其距离记作 d。函数 \(f: X \to C\) 称为 Lipschitz 函数，如果存在实数 \(c \geqslant 0\)，使得对 X 中所有 \(x\) 与 \(y\) 有
\end{enumerate}

\[
| f (x) - f (y) | \leqslant c d (x, y)
\]

（参见 FVR，IV，第 7 页，定义 1）。我们用 \(\mathrm{Lip}_b(\mathrm{X})\) 表示 X 上有界 Lipschitz 函数所成的集合。

\begin{enumerate}
\def\labelenumi{\alph{enumi})}
\item
  集合 \(\mathrm{Lip}_b(\mathrm{X})\) 是 \(\mathcal{C}(\mathrm{X})\) 的一个子代数。
\item
  赋以范数
\end{enumerate}

\[
\| f\| = \sup_{x\in \mathrm{X}}|f(x)| + \sup_{\substack{x,y\in \mathrm{X}\\ x\neq y}}\frac{|f(x) - f(y)|}{d(x,y)},
\]

以及对合 \(f \mapsto \overline{f}\) 后，代数 \(\mathrm{Lip}_b(\mathrm{X})\) 是一个交换对合 Banach 代数。它一般不是星代数。

\begin{enumerate}
\def\labelenumi{\alph{enumi})}
\setcounter{enumi}{2}
\item
  对任意交换 Banach 代数 A，我们用 \(\mathsf{X}_{m}(\mathsf{A})\) 表示度量空间 \(\mathsf{X}(\mathsf{A})\)，它赋有由 A 的对偶 \(A'\) 的范数所诱导的距离。度量空间 \(\mathsf{X}_{m}(\mathsf{A})\) 是完备的，且直径 \(\leqslant 2\)。
\item
  设 A 为幺交换 Banach 代数。Gelfand 变换把 A 映到 \(\operatorname{Lip}_{b}(\mathsf{X}_{m}(\mathsf{A}))\) 中。若 A 是半单的，则它诱导从 A 到其像上的一个同胚。
\end{enumerate}

以下我们假设 X 是紧度量空间。记 \(A = \operatorname{Lip}_{b}(X)\)。

\begin{enumerate}
\def\labelenumi{\alph{enumi})}
\setcounter{enumi}{4}
\item
  映射 \(\varphi\) 把 \(x \in X\) 映到特征标 \(f \mapsto f(x)\)，它是从 \(X\) 到 \(X(A)\) 上的一个同胚。
\item
  设 \(d_{1}\) 为由同胚 \(\varphi\) 诱导的 X 上的距离。则 \(d_{1}\) 与 d 等价。
\end{enumerate}

\chapter{局部紧交换群}

在本章中，除非另有说明，字母 G 表示一个赋有 Haar 测度（通常记作 dx）的局部紧交换群；对 \(p \in [1, +\infty]\)，空间 \(L^{p}(G, dx)\) 将简记为 \(L^{p}(G)\)，其范数记作 \(f \mapsto \|f\|_{p}\)。我们把 \(L^{1}(G)\) 等同于 \(\mathcal{M}^{1}(G)\) 的一个子空间，等同映射为 \(f \mapsto f \cdot dx\)。回忆 Haar 测度的支集等于 G（INT，VII，§1，第 1 节，注 3）；特别地（INT，III，§2，第 2 节，命题 9），从空间 \(\mathcal{K}(G)\) 到 \(L^{p}(G)\) 的典范映射对 \(p \in [1, +\infty]\) 是单射。对 \(p \neq +\infty\)，我们把 \(L^{p}(G)\) 等同于空间 \(\widetilde{\mathcal{F}}(G; \mathbf{C})\) 的一个子空间，后者是 G 上几乎处处有定义且取有限值的复值函数的类所成的空间（INT，IV，§3，第 5、6 节）。特别地，记号 \(L^{1}(G) \cap L^{2}(G)\) 表示 \(L^{1}(G)\) 与 \(L^{2}(G)\) 在该空间中的交。我们用 \(f \mapsto \widetilde{f}\) 表示对合代数 \(L^{1}(G)\) 上的对合（I，第 99 页，例 4）；\(\widetilde{f}(x) = \overline{f(x^{-1})}\) 对 G 中一切 \(x\) 成立。

回忆（INT，V，§5，第 3 节，定理 1）：若 μ 是局部紧拓扑空间 X 上的一个复测度，f 是 X 上一个局部 μ-可积函数，则以 f 为密度关于 μ 的测度 ν 记作 f·μ 或简记作 fμ。从 X 到 C 的函数 g 关于测度 ν 本质可积当且仅当 gf 关于 \(\mu\) 本质可积；此时有

\[
(f \cdot \mu) (g) = \int_ {\mathrm{X}} g d (f \cdot \mu) = \int_ {\mathrm{X}} g f d \mu = \mu (g f).
\]

若 G 是紧拓扑群，G 上的规范化 Haar 测度是 G 上唯一的 Haar 测度 \(\mu\)，它满足 \(\mu(\mathrm{G}) = 1\)。

若 \(X\) 是离散拓扑空间，\(X\) 上的计数测度是离散测度 \(\mu\)（在 \(X\) 上），它满足 \(\mu(\{x\}) = 1\) 对所有 \(x \in X\) 成立。若 \(G\) 是离散拓扑群，则 \(G\) 上的计数测度是 \(G\) 上的一个 Haar 测度。

我们还将用到下面两个引理。

引理 1.　设 X 为局部紧拓扑空间，\(\mu\) 为 X 上的一个正测度。设 \(x\) 为 \(\mu\) 的支集中的一个元素。设 U 为 \(x\) 的一个开邻域。存在一个支集含于 U 的函数 \(f \in \mathcal{K}_{+}(\mathrm{X})\)，使得 \(\int fd\mu = 1\)。

由测度支集的定义（INT，III，§2，第 2 节，定义 1），存在一个支集含于 U 的函数 \(g \in \mathcal{K}(\mathrm{X})\)，使得 \(\mu(g) \neq 0\)。我们有 \(\mu(|g|) > 0\)，因为 \(\mu\) 是正测度，而函数 \(f = \mu(|g|)^{-1}|g|\) 即具有所需的性质。

引理 2.　设 G 与 H 为分别以 \(e_{\mathrm{G}}\) 与 \(e_{\mathrm{H}}\) 为单位元的拓扑群。假设拓扑群 G 是 Hausdorff 的。设 f 为从 G 到 H 的一个拓扑群态射。假设对 G 中 \(e_{\mathrm{G}}\) 的每个邻域 U，存在 H 中 \(e_{\mathrm{H}}\) 的一个邻域 W，使得 \(f^{-1}(\mathrm{W}) \subset \mathrm{U}\)。则态射 f 是单的且是严格的（TG，III，第 16 页，定义 1）。

这些假设蕴含 \(\operatorname{Ker}(f)\) 包含于 G 中 \(e_{G}\) 的每个邻域，从而由于 G 是分离的，有 \(\operatorname{Ker}(f) = \{e_{\mathrm{G}}\}\)。于是 f 通过过渡到子空间所诱导的从 G 到 \(f(\mathrm{G})\) 的同态是双射；用 \(g: f(\mathrm{G}) \to \mathrm{G}\) 表示逆同态。这些假设进而蕴含 g 在 \(e_{H}\) 处连续，从而连续（TG，III，第 15 页，命题 23）。

\section{FOURIER 变换}

\subsection{局部紧交换群的酉特征标}

定义 1.　G 的酉特征标是从 G 到模为 1 的复数所成的乘法群 U 中的一个连续同态。

换句话说，酉特征标是 G 上的一个连续复值函数 \(\chi\)，满足：

\[
\chi (x y) = \chi (x) \chi (y), \qquad | \chi (x) | = 1 \qquad (x, y \in \mathrm{G}).
\]

在本章中，我们常把"酉特征标"简称为"特征标"。

设 E 为一维 Hilbert 空间，\(\chi\) 为 G 的一个酉特征标。把 \(x \in \mathrm{G}\) 映到 E 中比为 \(\chi(x)\) 的位似映射的映射是 G 在 E 上的一个连续线性等距表示。反之，G 在 E 上的每个有界连续线性表示都可由此过程得到，特别地它是酉的。

两个酉特征标的乘积、一个酉特征标的逆以及等于 1 的常函数都是酉特征标，这是直接的。因此，G 的酉特征标所成的集合 \(\widehat{\mathbf{G}}\) 对乘法构成一个群。这个群是交换的。另一方面，映射 \((\chi_1,\chi_2)\mapsto \chi_1\chi_2^{-1} = \chi_1\overline{\chi}_2\) 对紧开拓扑是连续的，从而赋有紧收敛拓扑的 \(\widehat{\mathbf{G}}\) 是一个拓扑群（TG，X，第 6 页，推论 2 与注 1）。

定义 2.　拓扑群 \(\widehat{\mathbf{G}}\) 称为 G 的对偶群。

由于 G 是局部紧的，映射 \((x,\chi)\mapsto\chi(x)\) 在 \(G\times\widehat{G}\) 上连续（TG，X，第 28 页，定理 3）。

回忆 \(\mathcal{M}^{1}(G)\) 表示 G 上有界复测度所成的幺对合 Banach 代数（I，第 99 页，例 4）。对每个有界复测度 \(\mu\in\mathcal{M}^{1}(G)\) 与每个 \(\chi\in\widehat{G}\)，我们记

\[
\chi (\mu) = \int_ {\mathrm{G}} \chi (x) d \mu (x)\tag{1}
\]

（参见 INT，VIII，§2，第 6 节）。

引理 1.　对每个 \(\chi \in \widehat{\mathbf{G}}\)，映射 \(\mu \mapsto \chi(\mu)\) 是对合 Banach 代数 \(\mathcal{M}^1(\mathbf{G})\) 的一个 Hermite 特征标。

由 INT，VIII，§3，第 3 节，命题 11，映射 \(\mu \mapsto \chi(\mu)\) 是对合 Banach 代数 \(\mathcal{M}^{1}(G)\) 的一个特征标。此外，我们有：

\[
\chi (\mu^ {*}) = \int_ {\mathrm{G}} \chi (x ^ {- 1}) d \overline {{\mu}} (x) = \int_ {\mathrm{G}} \overline {{\chi (x)}} d \overline {{\mu}} (x) = \overline {{\chi (\mu)}}
\]

因此这个特征标是 Hermite 的。

这样我们定义了一个从 \(\widehat{G}\) 到 \(\mathsf{X}(\mathcal{M}^{1}(\mathrm{G}))\) 中的映射；它将称为典范映射。

设 \(\chi \in \widehat{\mathbf{G}}\)。\(\mu \mapsto \chi(\mu)\) 在 \(\mathrm{L}^1(\mathrm{G})\) 上的限制非零（参见 INT，VIII，§2，第 7 节，命题 10）。因此，通过限制到对合 Banach 子代数 \(\mathrm{L}^1(\mathrm{G})\)，我们得到一个从 \(\widehat{\mathbf{G}}\) 到 \(\mathsf{X}(\mathrm{L}^1(\mathrm{G}))\) 中的映射，称为典范映射。它把 \(\chi \in \widehat{\mathbf{G}}\) 对应到 Hermite 特征标

\[
f \mapsto \chi (f) = \chi (f \cdot d x) = \int_ {\mathrm{G}} f (x) \chi (x) d x \qquad (f \in \mathrm{L} ^ {1} (\mathrm{G}))\tag{2}
\]

即 \(L^{1}(G)\) 的特征标。

命题 1.　从 \(\widehat{\mathbf{G}}\) 到 \(\mathsf{X}(\mathrm{L}^{1}(\mathbf{G}))\) 的典范映射是一个同胚。

把这个典范映射记作 ev，并对 \(\chi \in \widehat{\mathbf{G}}\)，用 \(\mathrm{ev}_{\chi}\) 表示特征标 \(\chi\) 在 ev 下的像，即 Hermite 特征标 \(f \mapsto \chi(f)\)，它定义在 \(\mathrm{L}^1(\mathrm{G})\) 上。把 ev 视为从 \(\widehat{\mathbf{G}}\) 到 \(\mathrm{L}^1(\mathrm{G})\) 的对偶中的映射时，它是 \(\widehat{\mathbf{G}}\) 到 \(\mathrm{L}^{\infty}(\mathrm{G})\)（赋有弱拓扑 \(\sigma(\mathrm{L}^{\infty}(\mathrm{G}), \mathrm{L}^{1}(\mathrm{G}))\)）中的单射与 \(\mathrm{L}^{\infty}(\mathrm{G})\) 到赋有点态收敛拓扑的 \(\mathrm{L}^1(\mathrm{G})\) 的对偶中的单射的复合。由于 \(\widehat{\mathbf{G}}\) 是 \(\mathrm{L}^{\infty}(\mathrm{G})\) 的有界子集，由 Lebesgue 定理（INT，IV，§3，第 7 节，定理 6），第一个映射连续。第二个映射按定义也是连续的。这就证明了映射 ev 连续。

若 \(\chi \in \widehat{\mathrm{G}}\) 且 \(f \in \mathrm{L}^1(\mathrm{G})\)，则有 \(\mathrm{ev}_{\chi}(\varepsilon_x * f) = \chi(x)\mathrm{ev}_{\chi}(f)\)。取 \(f\) 使 \(\mathrm{ev}_{\chi}(f) \neq 0\)，便推出映射 ev 是单射。

设 \(\zeta\in\mathsf{X}(\mathrm{L}^{1}(\mathrm{G}))\)，并设 \(f\in\mathrm{L}^{1}(\mathrm{G})\) 满足 \(\zeta(f)\neq0\)。我们用如下方式定义映射 \(\chi\colon\mathrm{G}\to\mathbf{C}\)：对 \(x\in\mathrm{G}\)，令

\[
\chi (x) = \frac {\zeta (\varepsilon_ {x} * f)}{\zeta (f)}.\tag{3}
\]

来定义它。我们有 \(\chi(e) = 1\)。由于映射 \(x \mapsto \varepsilon_x * f = \gamma(x)(f)\)（从 G 到 L \(^1\) (G)）连续（INT，VIII，§2，第 5 节，命题 8），映射 \(\chi\) 连续。它是有界的，因为对所有 \(x \in \mathbf{G}\) 有

\[
| \chi (x) | \leqslant \frac {\| \varepsilon_ {x} * f \|}{| \zeta (f) |} = \frac {\| f \|}{| \zeta (f) |}
\]

（I，第 29 页，定理 1 及 INT，VIII，同前）。

现设 \(\mathfrak{B}\) 为由紧邻域组成的 \(e\) 的邻域滤子的一个基。对每个 \(V \in \mathfrak{B}\)，设 \(g_V\) 为一个在 \(V\) 外为零且积分等于 1 的正连续函数（II，第 200 页，引理 1）。于是对每个函数 \(h \in L^1(G)\) 有

\[
\varepsilon_ {x} * h = \lim _ {\mathrm{V}, \mathfrak {B}} (\varepsilon_ {x} * h) * g _ {\mathrm{V}} = \lim _ {\mathrm{V}, \mathfrak {B}} (\varepsilon_ {x} * g _ {\mathrm{V}} * h),
\]

在 L \(^{1}\) (G) 中成立，其中极限是沿 \(\mathfrak{B}\) 的截口滤子取的（INT，VIII，§4，第 7 节，命题 20）。特别地，由于 \(\zeta(\varepsilon_x * g_V * f) = \zeta(\varepsilon_x * g_V)\zeta(f)\)，由此推出

\[
\chi (x) = \lim _ {\mathrm{V}, \mathfrak {B}} \zeta (\varepsilon_ {x} * g _ {\mathrm{V}}).
\]

对每个 \(h\in \mathrm{L}^1 (\mathrm{G})\)，我们得到

\[
\zeta (\varepsilon_ {x} * h) = \lim _ {\mathrm{V}, \mathfrak {B}} \zeta (\varepsilon_ {x} * g _ {\mathrm{V}} * h) = \zeta (h) \lim _ {\mathrm{V}, \mathfrak {B}} \zeta (\varepsilon_ {x} * g _ {\mathrm{V}}) = \chi (x) \zeta (h).
\]

从而，对 \(x, y \in G\) 有：

\[
\chi (x y) = \frac {\zeta (\varepsilon_ {x} * \varepsilon_ {y} * f)}{\zeta (f)} = \frac {\chi (x) \zeta (\varepsilon_ {y} * f)}{\zeta (f)} = \chi (x) \chi (y),
\]

这就证明了 \(\chi\) 是从 G 到 \(\mathbf{C}^*\) 的一个同态。由于 \(\chi\) 有界且连续，它是 G 的一个酉特征标。此外，若 \(g \in \mathrm{L}^{1}(\mathrm{G})\)，则有

\[
g * f = \int_ {\mathrm{G}} (\varepsilon_ {x} * f) g (x) d x
\]

在 L \(^{1}\) (G) 中成立（INT，VIII，§1，第 5 节，命题 7），由此得

\[
\begin{array}{c} \zeta (g) \zeta (f) = \zeta (g * f) = \int_ {\mathrm{G}} \zeta (\varepsilon_ {x} * f) g (x) d x \\ = \zeta (f) \int_ {\mathrm{G}} \chi (x) g (x) d x = \mathrm{ev} _ {\chi} (g) \zeta (f) \end{array}
\]

（INT，VI，§1，第 1 节，命题 1），这表明 \(\zeta = \mathrm{ev}_{\chi}\)。因此 ev 是满的，从而是双射。

最后证明逆映射 \(ev^{-1}\) 连续。设 \(\zeta \in \mathsf{X}(\mathrm{L}^{1}(\mathrm{G}))\)。设 \(f \in \mathrm{L}^{1}(\mathrm{G})\) 为使 \(\zeta(f) \neq 0\) 的函数。由 \(\xi \in \mathsf{X}(\mathrm{L}^{1}(\mathrm{G}))\)（满足 \(\xi(f) \neq 0\)）所成的集合 W 是 \(\zeta\) 在 \(\mathsf{X}(\mathrm{L}^{1}(\mathrm{G}))\) 中的一个开邻域。对每个 \(\xi \in W\)，前文的论证表明 \(\mathrm{ev}^{-1}(\xi)\) 是特征标

\[
x \mapsto \frac {\xi (\varepsilon_ {x} * f)}{\xi (f)}.
\]

设 \(\mathfrak{F}\) 为 \(W \subset X(L^{1}(G))\) 上收敛于 \(\zeta\) 的一个滤子。由于集合 \(X(L^{1}(G))\) 在 \(L^{\infty}(G)\) 中有界从而等度连续，单收敛的一致结构与紧收敛的一致结构一致（TG，X，第 16 页，定理 1）。设 K 为 G 的一个紧子集。诸 \(\varepsilon_{x} * f\)（\(x \in K\)）所成的集合在 \(L^{1}(G)\) 中是紧的（INT，VIII，§2，第 5 节，命题 8）。于是我们有

\[
\lim _ {\xi , \mathfrak {F}} \xi (\varepsilon_ {x} * f) = \zeta (\varepsilon_ {x} * f)
\]

对 \(x \in K\) 一致地成立。由此推出

\[
\lim _ {\xi , \mathfrak {F}} \mathrm{ev} ^ {- 1} (\xi) = \mathrm{ev} ^ {- 1} (\zeta),
\]

即 ev \(^{-1}\) 在 \(\zeta\) 处连续。命题证毕。

从现在起，我们把 G 的酉特征标 \(\chi\) 与特征标 \(f\mapsto\int_{\mathrm{G}}f(x)\chi(x)dx\)（\(\mathrm{L}^{1}(\mathrm{G})\) 上的）等同起来。

注.　1) *命题 1 中从 \(\widehat{G}\) 到 \(X(L^{1}(G))\) 的映射的双射性，是局部紧群 H（不必交换）的连续表示与代数 \(L^{1}(H)\) 的连续表示之间对应关系的一个特殊情形。*

\begin{enumerate}
\def\labelenumi{\arabic{enumi})}
\setcounter{enumi}{1}
\tightlist
\item
  从 \(\widehat{G}\) 到 \(\mathsf{X}(\mathcal{M}^{1}(\mathsf{G}))\) 的典范映射一般不是满射（II，第 308 页，习题 14）。
\end{enumerate}

推论 1.　L \(^{1}\) (G) 的每个特征标都是 Hermite 的。从 X(Stell(G))（I，第 125 页，定义 9）到 X(L \(^{1}\) (G)) 的典范映射是一个同胚。

第一个断言由命题 1 与引理 1 限制到 \(L^{1}(G)\) 得到。第二个断言由第一个断言以及 I，第 124 页，命题 20 的推论得到。

推论 2.　拓扑群 \(\widehat{\mathbf{G}}\) 是局部紧的。

事实上，\(\mathsf{X}(\mathrm{L}^{1}(\mathrm{G}))\) 是局部紧的（I，第 29 页，定理 1 的推论）。

我们将把 \(\widehat{G}\) 与 \(\mathsf{X}(\mathrm{L}^1 (\mathrm{G}))\) 以及 \(\mathsf{X}(\mathrm{Stell}(\mathrm{G}))\) 等同起来。对 \(x\in G\) 与 \(\chi \in \widehat{G}\)，我们用 \(\langle \chi ,x\rangle\) 表示复数 \(\chi (x)\)，它属于 \(\mathbf{U}\)。

我们称 \(x\) 与 \(\chi\) 正交，如果 \(\langle \chi, x \rangle = 1\)。设 A 为 G（相应地，\(\widehat{\mathbf{G}}\)）的一个子集；\(\widehat{\mathbf{G}}\)（相应地，G）中与 A 正交的元素所成的集合是 \(\widehat{\mathbf{G}}\)（相应地，G）的一个闭子群，称为 A 的正交子群，记作 \(\mathrm{A}^{\perp}\)。G 的正交子群仅由 \(e\) 组成。

对 \(x \in G\)，设 \(\eta(x)\) 为从 \(\widehat{G}\) 到 U 的、由 \(\chi \mapsto \langle \chi, x \rangle\) 定义的映射。由 \(\widehat{G}\) 中乘法的定义，映射 \(\eta(x)\) 是一个群同态。由于映射 \((x, \chi) \mapsto \langle \chi, x \rangle\)（从 \(G \times \widehat{G}\) 到 U）连续（TG，X，第 28 页，定理 3），它连续。这样我们定义了一个映射 \(\eta\)，称为典范映射，它从 G 映到二重对偶群 \(\widehat{\widehat{G}}\) 中；它是一个群同态。此外，映射 \(\eta\) 连续（TG，X，第 28 页，定理 3）。我们稍后将证明（II，第 220 页，定理 2）\(\eta\) 是从 G 到 \(\widehat{\widehat{G}}\) 上的拓扑群同构。

设 G 与 H 为局部紧交换群，\(\varphi\colon\mathrm{G}\to\mathrm{H}\) 为一个拓扑群态射。对每个 \(\chi\in\widehat{\mathrm{H}}\)，映射 \(\chi\circ\varphi\) 是 G 的一个特征标，记作 \(\widehat{\varphi}(\chi)\)。这一定义可以写成公式

\[
\langle \chi , \varphi (x) \rangle = \langle \widehat {\varphi} (\chi), x \rangle\tag{4}
\]

对所有 \(\chi\in\widehat{H}\) 与 \(x\in G\) 成立。由此推出 \(\widehat{\varphi}\) 是从拓扑群 \(\widehat{H}\) 到拓扑群 \(\widehat{G}\) 的一个态射；我们称 \(\widehat{\varphi}\) 为态射 \(\varphi\) 的对偶。

设 K 为局部紧交换群，\(\psi\colon\mathrm{H}\to\mathrm{K}\) 为一个拓扑群态射。定义表明 \(\widehat{\psi\circ\varphi}=\widehat{\varphi}\circ\widehat{\psi}\)。若 \(\varphi\) 是 G 的恒等映射，则 \(\widehat{\varphi}\) 是 \(\widehat{\mathrm{G}}\) 的恒等映射。特别地，若 \(\varphi\) 是拓扑群同构，则 \(\widehat{\varphi}\) 也是，且 \(\widehat{\varphi}^{-1}\) 是 \(\varphi^{-1}\) 的对偶。

引理 2.　设 G 与 H 为局部紧交换群，\(f\colon\mathrm{H}\to\mathrm{G}\) 为一个拓扑群态射。\(\widehat{f}\) 的核是 \(f\) 的像的正交子群。

按定义，\(\chi \in \mathrm{Ker}(\widehat{f})\) 当且仅当 \(\chi\) 在 \(f\) 的像上的限制是平凡的。

命题 2.　设 \(n \geqslant 0\) 为整数，\(\mathrm{G}_{1}, \ldots, \mathrm{G}_{n}\) 为局部紧交换群。设 G 为诸群 \(\mathrm{G}_{j}\)（\(1 \leqslant j \leqslant n\)）的乘积群。对 \(1 \leqslant j \leqslant n\)，设 \(\lambda_{j}\) 为 \(\mathrm{G}_{j}\) 到 G 中的单射，它把 \(x \in \mathrm{G}_{j}\) 映到元素 \((x_{k})\)（\(x_{k} = e\) 若 \(k \neq j\)，而 \(x_{j} = x\)）。映射

\[
(\widehat {\lambda} _ {j}) _ {1 \leqslant j \leqslant n} \colon \widehat {\mathrm{G}} \to \prod_ {1 \leqslant j \leqslant n} \widehat {\mathrm{G}} _ {j}
\]

是一个拓扑群同构。

设 \(m\) 为从 \(\widehat{\mathbf{G}}^n\) 到 \(\widehat{\mathbf{G}}\) 的乘积映射，并对每个 \(j\)（\(1 \leqslant j \leqslant n\)），设 \(\pi_j\) 为 \(\mathbf{G}\) 到 \(\mathbf{G}_j\) 上的投影。设 \(\mu\) 为拓扑群态射 \(m \circ (\widehat{\pi}_j)_j\)，它从 \(\prod \widehat{\mathbf{G}}_j\) 到 \(\widehat{\mathbf{G}}\)，于是

\[
\langle \mu ((\chi_ {j})), (x _ {j}) \rangle = \prod_ {j = 1} ^ {n} \langle \chi_ {j}, x _ {j} \rangle .
\]

映射 \(\mu\) 连续，并且可以验证 \(\mu\) 与 \((\widehat{\lambda}_j)\) 互为逆双射。命题得证。

注.　紧交换群的无限乘积的对偶群的计算是下文 II，第 234 页，推论 4 的主题。至于作为局部紧群的任意乘积的局部紧交换群的情形，则可由这两个命题得到，因为在这种乘积中除有限多个因子外都是紧的（TG，I，第 66 页，命题 14，b））。

\subsection{Fourier 变换的定义}

定义 3.　设 \(\mu \in \mathcal{M}^{1}(\mathrm{G})\) 为 G 上的一个有界复测度。\(\mu\) 的 Fourier 变换是函数 \(\mathcal{F}_{\mathrm{G}}(\mu)\)，它定义在 \(\widehat{\mathrm{G}}\) 上，由下式给出：

\[
\mathcal {F} _ {\mathrm{G}} (\mu) (\widehat {x}) = \int_ {\mathrm{G}} \overline {{\langle \widehat {x} , x \rangle}} d \mu (x).\tag{5}
\]

我们把 \(\mu\) 的 Fourier 余变换称为函数 \(\overline{\mathcal{F}}_{\mathrm{G}}(\mu)\)，它定义在 \(\widehat{\mathrm{G}}\) 上，由下式给出：

\[
\overline {{\mathcal {F}}} _ {\mathrm{G}} (\mu) (\widehat {x}) = \int_ {\mathrm{G}} \langle \widehat {x}, x \rangle d \mu (x).\tag{6}
\]

当所考虑的群 G 没有歧义时，我们也写作 \(\mathcal{F}(\mu)\) 与 \(\overline{\mathcal{F}}(\mu)\)。我们有时也记 \(\widehat{\mu} = \mathcal{F}_{\mathrm{G}}(\mu)\)。

命题 3.　对每个测度 \(\mu \in \mathcal{M}^{1}(\mathrm{G})\)，函数 \(\mathcal{F}_{\mathrm{G}}(\mu)\) 与 \(\overline{\mathcal{F}}_{\mathrm{G}}(\mu)\) 连续且有界。映射 \(\mathcal{F}_{\mathrm{G}}: \mu \mapsto \mathcal{F}_{\mathrm{G}}(\mu)\) 与 \(\overline{\mathcal{F}}_{\mathrm{G}}: \mu \mapsto \overline{\mathcal{F}}_{\mathrm{G}}(\mu)\) 是从对合代数 \(\mathcal{M}^{1}(\mathrm{G})\) 到 \(\widehat{\mathrm{G}}\) 上有界连续函数所成的对合代数（I，第 99 页，例 1）中的连续态射。

设 \(\mu \in \mathcal{M}^1 (\mathrm{G})\)。对每个 \(\chi \in \widehat{\mathrm{G}}\) 有

\[
| \mathcal {F} (\mu) (\chi) | = \left| \int_ {\mathrm{G}} \overline {{\langle \chi , x \rangle}} d \mu (x) \right| \leqslant \| \mu \| _ {1},\tag{7}
\]

因此 \(\mu\) 的 Fourier 变换是有界的。类似地，可以验证 \(\overline{\mathcal{F}}(\mu)\) 是有界的。

若 \(\chi\) 趋于 \(\chi_0\)（于 \(\widehat{\mathbf{G}}\) 中），则函数 \(\chi\)（在 \(\mathbf{G}\) 上）在每个紧集上一致趋于 \(\chi_0\)，同时保持被常函数 1 控制，而 1 属于 \(\mathrm{L}^1 (\mathrm{G},\mu)\)。由 Lebesgue 定理（INT，IV，§3，第 7 节，定理 6），由此得 \(\mathcal{F}(\mu)(\chi)\) 趋于 \(\mathcal{F}(\mu)(\chi_0)\)。因此 \(\mathcal{F}(\mu)\) 连续。对 \(\overline{\mathcal{F}} (\mu)\) 同样成立。

对每个 \(\chi \in \widehat{\mathbf{G}}\)，映射 \(\mu \mapsto \chi(\mu) = \int \langle \chi, x \rangle d\mu(x)\) 是 \(\mathcal{M}^1(\mathbf{G})\) 的一个 Hermite 特征标（II，第 202 页，引理 1）。这蕴含 \(\mathcal{F}\) 与 \(\overline{\mathcal{F}}\) 是从 \(\mathcal{M}^1(\mathbf{G})\) 到 \(\widehat{\mathbf{G}}\) 上有界连续函数所成的对合代数中的对合代数态射。不等式 (7) 表明这些态射连续。

G 的 Fourier 变换（相应地，G 的 Fourier 余变换）是映射 \(\mu \mapsto \mathcal{F}_{\mathrm{G}}(\mu)\)（相应地，映射 \(\mu \mapsto \overline{\mathcal{F}}_{\mathrm{G}}(\mu)\)），它从 \(\mathcal{M}^1 (\mathrm{G})\) 到 \(\mathcal{C}_b(\widehat{\mathrm{G}})\) 中。

我们记下一些有用的公式，对 \(\mu\in\mathcal{M}^{1}(G)\)、\(x\in G\) 与 \(\chi\in\widehat{G}\)：

\begin{enumerate}
\def\labelenumi{(\arabic{enumi})}
\setcounter{enumi}{7}
\tightlist
\item
\end{enumerate}

\[
\overline {{\mathcal {F}}} (\mu) (\chi) = \mathcal {F} (\mu) (\chi^ {- 1}) = \overline {{\mathcal {F} (\overline {{\mu}}) (\chi)}},\tag{9}
\]

\[
\| \mathcal {F} (\mu) \| _ {\infty} = \| \overline {{\mathcal {F}}} (\mu) \| _ {\infty} \leqslant \| \mu \| _ {1},\tag{10}
\]

\[
\left\{ \begin{array}{l} \mathcal {F} (\varepsilon_ {x}) (\chi) = \overline {{\langle \chi , x \rangle}}, \\ \overline {{\mathcal {F}}} (\varepsilon_ {x}) (\chi) = \langle \chi , x \rangle , \end{array} \right.
\]

（特别地 \(\mathcal{F}(\varepsilon_e) = \overline{\mathcal{F}} (\varepsilon_e) = 1\)），

\begin{enumerate}
\def\labelenumi{(\arabic{enumi})}
\setcounter{enumi}{10}
\tightlist
\item
\end{enumerate}

\[
\left\{ \begin{array}{l} \mathcal {F} (\varepsilon_ {x} * \mu) (\chi) = \overline {{\langle \chi , x \rangle}} \mathcal {F} (\mu) (\chi), \\ \overline {{\mathcal {F}}} (\varepsilon_ {x} * \mu) (\chi) = \langle \chi , x \rangle \overline {{\mathcal {F}}} (\mu) (\chi), \end{array} \right.\tag{12}
\]

\[
\left\{ \begin{array}{l} \mathcal {F} (\chi \cdot \mu) = \varepsilon_ {\chi} * \mathcal {F} (\mu), \\ \overline {{\mathcal {F}}} (\chi \cdot \mu) = \varepsilon_ {\chi^ {- 1}} * \overline {{\mathcal {F}}} (\mu). \end{array} \right.
\]

公式 (8)、(9)、(10) 与 (11) 由定义直接得出。我们来证明公式 (12) 中的第一个，第二个是类似的。对一切 \(\xi\)（于 \(\widehat{G}\) 中），等式

\[
\begin{array}{c} \mathcal {F} (\chi \cdot \mu) (\xi) = \int_ {\mathrm{G}} \overline {{\langle \xi , x \rangle}} \langle \chi , x \rangle d \mu (x) = \int_ {\mathrm{G}} \overline {{\langle \xi \chi^ {- 1} , x \rangle}} d \mu (x) \\ = \mathcal {F} (\mu) (\xi \chi^ {- 1}) = (\varepsilon_ {\chi} * \mathcal {F} (\mu)) (\xi). \end{array}
\]

此外注意，对所有 \(\chi\in\widehat{\mathrm{G}}\) 以及所有测度 \(\mu\) 与 \(\nu\)（在 \(\mathcal{M}^{1}(\mathrm{G})\) 中），有

\[
(\varepsilon_ {\chi} * \mathcal {F} (\mu)) (\varepsilon_ {\chi} * \mathcal {F} (\nu)) = \varepsilon_ {\chi} * (\mathcal {F} (\mu) \mathcal {F} (\nu)),\tag{13}
\]

因为这个等式两边都是 \(\widehat{G}\) 上的函数，其在 \(\xi\in\widehat{G}\) 处的值为

\[
\mathcal {F} (\mu) (\xi \chi^ {- 1}) \mathcal {F} (\nu) (\xi \chi^ {- 1}).
\]

设 H 为局部紧交换群，\(\varphi\colon\mathrm{G}\to\mathrm{H}\) 为一个连续态射。设 \(\mu\in\mathcal{M}^{1}(\mathrm{G})\)。像测度 \(\varphi(\mu)\) 已有定义（INT，V，§6，第 1 节，注 1），并且由此可得 \(\mathcal{F}_{\mathrm{H}}(\varphi(\mu))=\mathcal{F}_{\mathrm{G}}(\mu)\circ\widehat{\varphi}\)（参见 INT，V，§6，第 4 节，命题 7）。

通过限制到子代数 \(L^{1}(G)\)（\(\mathcal{M}^{1}(G)\) 的），我们得到 \(L^{1}(G)\) 上 Fourier 变换与 Fourier 余变换的定义。因此，对 \(f \in L^{1}(G)\) 与 \(\chi \in \widehat{G}\) 有：

\[
\mathcal {F} _ {\mathrm{G}} (f) (\chi) = \int_ {\mathrm{G}} \overline {{\langle \chi , x \rangle}} f (x) d x, \quad \overline {{\mathcal {F}}} _ {\mathrm{G}} (f) (\chi) = \int_ {\mathrm{G}} \langle \chi , x \rangle f (x) d x.\tag{14}
\]

特别地，\(\mathcal{F}_{\mathrm{G}}(f) = \overline{\overline{\mathcal{F}}_{\mathrm{G}}(\overline{f})}\)。我们还有

\[
\overline {{\mathcal {F}}} (f) (\chi) = \chi (f)\tag{15}
\]

对所有 \(f \in \mathrm{L}^{1}(\mathrm{G})\) 与每个 \(\chi \in \widehat{G}\) 成立。

设 \(\sigma\) 为 G 的一个自同构，\(\Delta\) 为 \(\sigma\) 的模（INT，VII，§1，第 4 节，定义 4）。对 \(f \in \mathrm{L}^1(\mathrm{G})\) 有

\[
\mathcal {F} (f \circ \sigma) = \Delta^ {- 1} \mathcal {F} (f) \circ \widehat {\sigma} ^ {- 1}\tag{16}
\]

（参见同前，公式 (31)）。

若把 \(\widehat{G}\) 与 \(X(L^{1}(G))\) 等同起来（II，第 202 页，命题 1），则 Fourier 余变换不是别的，正是 Banach 代数 \(L^{1}(G)\) 的 Gelfand 变换（I，第 7 页，定义 5）。

命题 4.　Fourier 变换与 Fourier 余变换都是从 L \(^{1}\) (G) 到代数 \(\mathcal{C}_{0}(\widehat{\mathrm{G}})\)（\(\widehat{\mathrm{G}}\) 上在无穷远处为 0 的连续函数所成的代数）中的单射对合代数态射。

Fourier 余变换是从 \(L^{1}(G)\) 到 \(\widehat{G}\) 上有界连续函数代数中的一个对合代数态射。由于它与 Gelfand 变换等同，其像含于 \(\mathcal{C}_{0}(\widehat{\mathrm{G}})\)（I，第 37 页，命题 5），其核是 \(L^{1}(G)\) 的根（I，第 38 页，命题 8），而根为零（I，第 126 页，命题 22 的推论）。

我们稍后将看到（II，第 221 页，命题 13 的推论），\(\mathcal{M}^{1}(\mathrm{G})\) 上的 Fourier 余变换也是单射的。

注.　空间 \(L^{1}(G)\) 上的 Fourier 变换依赖于 Haar 测度 dx 的选取，这与 \(\mathcal{M}^{1}(G)\) 上的 Fourier 变换不同。若把 dx 换成测度 \(a \cdot dx\)（其中 \(a > 0\)），则对 G 上任何可积函数 f，f 的 Fourier 变换是 \(a\widehat{f}\)，其中 \(\widehat{f}\) 是关于测度 dx 定义的 Fourier 变换。

考虑群 G 的星代数 Stell(G)（I，第 125 页，定义 9），并把 L \(^{1}\) (G) 等同于 Stell(G) 的一个稠密子代数（I，第 126 页，命题 22）。

命题 5.　由连续性，Fourier 变换与 Fourier 余变换都唯一地扩张为从 Stell(G) 到 \(\mathcal{C}_{0}(\widehat{\mathrm{G}})\) 上的星代数同构。

Fourier 余变换由连续性扩张为从 Stell(G) 到 \(\mathcal{C}_{0}(\widehat{\mathrm{G}})\) 中的一个星代数态射。若把 \(\widehat{G}\) 与 \(\mathrm{X}(\mathrm{Stell}(\mathrm{G}))\) 等同起来（II，第 204 页，推论 1 与 II，第 202 页，命题 1），这个扩张就是 Stell(G) 的 Gelfand 变换。由 I，第 108 页，定理 1，它是一个同构。关于 Fourier 变换的断言由此得出。

我们总是用 \(\overline{\mathcal{F}}\) 与 \(\mathcal{F}\) 表示命题 5 中的同构。

推论.　\(L\)，即 \(\mathrm{L}^{1}(\mathrm{G})\) 在 G 的 Fourier 变换下的像，在 \(\mathcal{C}_{0}(\widehat{\mathrm{G}})\) 中稠密。

由于 \(L^{1}(G)\) 在 \(\text{Stell}(G)\) 中稠密，这由命题 5 得出。

命题 6.　假设 G 是紧的。规范化 Haar 测度 dx 属于 \(\mathcal{M}^{1}(\mathrm{G})\)，其 Fourier 变换是 \(\varphi_{e}\)，即 \(\{e\}\) 的特征函数。

设 \(\chi \in \widehat{\mathrm{G}}\)。由于 \(\varepsilon_y * dx = dx\) 对所有 \(y \in \mathrm{G}\) 成立，有

\[
\mathcal {F} (d x) (\chi) = \overline {{\langle \chi , y \rangle}} \mathcal {F} (d x) (\chi)
\]

（公式 (11)）。若 \(\chi \neq 1\)，则存在 \(y \in \mathrm{G}\) 使 \(\langle \chi, y \rangle \neq 1\)，从而 \(\mathcal{F}(dx)(\chi) = 0\)。若 \(\chi = 1\)，则 \(\mathcal{F}(dx)(\chi) = \int_{\mathrm{G}} dx = 1\)，因为测度 \(dx\) 是规范化的。

\subsection{Plancherel 定理}

我们用 A(G) 表示 \(L^{1}(G)\) 的由诸函数 \(f * g\)（\(f, g \in L^{1}(G) \cap L^{2}(G)\)）生成的向量子空间。

命题 7.　空间 A(G) 是 L¹(G) 的一个自伴理想。它含于 L¹(G) ∩ L²(G)，也含于 C(G) 在 L¹(G) 中的像内。

设 \(f \in \mathrm{L}^1(\mathrm{G})\)。对每个 \(g \in \mathrm{L}^2(\mathrm{G})\)，有 \(f * g \in \mathrm{L}^1(\mathrm{G})\)。因此空间 \(\mathrm{L}^1(\mathrm{G}) \cap \mathrm{L}^2(\mathrm{G})\) 是 \(\mathrm{L}^1(\mathrm{G})\) 的一个理想，空间 \(\mathrm{A}(\mathrm{G})\) 亦然。理想 \(\mathrm{A}(\mathrm{G})\) 是自伴的。

设 \(f\) 与 \(g\) 属于 \(\mathrm{L}^2(\mathrm{G})\)。于是卷积 \(f * g\) 是由下式给出的连续函数的类

\[
y \mapsto \int_ {\mathrm{G}} f (y x ^ {- 1}) g (x) d x
\]

第二个断言由此得出。

由于 \(\chi(f*g) = (\chi f)*( \chi g)\) 对 \(\chi \in \widehat{\mathrm{G}}\)、\(f \in \mathrm{L}^1(\mathrm{G})\) 与 \(g \in \mathrm{L}^1(\mathrm{G})\) 成立（INT，VIII，§3，第 1 节，命题 6），我们有 \(\chi h \in \mathrm{A}(\mathrm{G})\) 对所有 \(h \in \mathrm{A}(\mathrm{G})\) 与 \(\chi \in \widehat{\mathrm{G}}\) 成立。由于 \(\varepsilon_x * f = \gamma(x)f\)，且线性表示 \(\gamma\) 在所有 \(\mathrm{L}^p(\mathrm{G})\)（对每个 \(p\)）上都是等距的（INT，VIII，§2，第 5 节，命题 8），我们有 \(\varepsilon_x * f \in \mathrm{A}(\mathrm{G})\) 对所有 \(x \in \mathrm{G}\) 与 \(f \in \mathrm{A}(\mathrm{G})\) 成立。

我们用 \(\widehat{\mathrm{A}}(\mathrm{G})\) 表示 \(\mathrm{A}(\mathrm{G})\) 在 Fourier 变换下的像。它是 \(\mathcal{C}_0(\widehat{\mathrm{G}})\) 的一个子空间。

命题 8.　存在一个滤子基 \(\mathfrak{B}\)（在 \(\mathrm{A}(\mathrm{G}) \cap \mathcal{K}_{+}(\mathrm{G})\) 上），使得下列条件满足：

\begin{enumerate}
\def\labelenumi{(\roman{enumi})}
\item
  对每个这样的 \(\varphi\)（\(\mathfrak{B}\) 的一个成员的元素），有 \(\|\varphi\|_1 = 1\) 且 \(\|\mathcal{F}(\varphi)\|_{\infty} \leqslant 1\)；
\item
  我们有
\end{enumerate}

\[
\lim _ {\varphi , \mathfrak {B}} \varphi \cdot d x = \varepsilon_ {e}
\]

在 G 上具有紧支集的测度所成的空间 \(\mathcal{C}'(\mathrm{G})\) 中成立，该空间赋有在 \(\mathcal{C}(\mathrm{G})\) 的紧子集上一致收敛的拓扑；

\begin{enumerate}
\def\labelenumi{(\roman{enumi})}
\setcounter{enumi}{2}
\tightlist
\item
  我们有
\end{enumerate}

\[
\lim _ {\varphi , \mathfrak {B}} \mathcal {F} (\varphi) = 1
\]

对 \(\widehat{\mathbf{G}}\) 上的紧收敛拓扑成立；

\begin{enumerate}
\def\labelenumi{(\roman{enumi})}
\setcounter{enumi}{3}
\tightlist
\item
  对 \(p = 1\) 或 \(p = 2\)，以及对每个 \(f \in \mathrm{L}^p(\mathrm{G})\)，有 \(\varphi * f \in \mathrm{A}(\mathrm{G})\) 对每个 \(\varphi\)（属于 \(\mathfrak{B}\) 的某个元素）成立，且
\end{enumerate}

\[
\lim _ {\varphi , \mathfrak {B}} \varphi * f = f
\]

在 L \(^{p}\) (G) 中成立。

设 \(K_{0}\) 为 G 中 e 的一个固定的紧邻域。设 \(B_{0}\) 为 G 中 e 的邻域滤子的一个基，由含于 \(K_{0}\) 的紧对称邻域组成（参见 TG，III，第 4 页）。对 \(K \in B_{0}\)，设 \(X_{K}^{\prime}\) 为诸函数 \(\psi \in \mathcal{K}_{+}(G)\)（满足 \(\operatorname{Supp}(\psi) \subset K\) 且 \(\int \psi(x)dx = 1\)）所成的集合；它非空（II，第 200 页，引理 1）。设 \(X_{K}\) 为诸函数 \(\psi * \psi\)（\(\psi \in X_{K}^{\prime}\)）所成的集合。它非空且含于 \(A(G) \cap \mathcal{K}_{+}(G)\)。集合 \(\mathfrak{B}\)（其元素为诸集合 \(X_{K}\)，其中 K 取遍 \(B_{0}\)）是 \(A(G) \cap \mathcal{K}_{+}(G)\) 上的一个滤子基。我们证明 \(\mathfrak{B}\) 满足所需的性质。

若 \(X \in \mathfrak{B}\) 且 \(\varphi \in X\)，则有 \(\| \varphi \|_1 = \int_{G} \varphi(x) dx = 1\)，从而 \(\| \mathcal{F}(\varphi) \|_\infty \leqslant 1\)，这就建立了性质 (i)。

性质 (ii) 由 INT，VIII，§2，第 7 节，引理 4 的推论 1 得出。\(\widehat{\mathbf{G}}\) 的紧子集是 \(\mathcal{C}(\mathbf{G})\) 的紧子集，因而 (ii) 蕴含 \(\lim_{\varphi, \mathfrak{B}} \mathcal{F}(\varphi) = 1\)（对 \(\widehat{\mathbf{G}}\) 上的紧收敛拓扑），即 (iii)。

最后，设 \(p=1\) 或 \(p=2\)。设 \(f\in\mathrm{L}^{p}(\mathrm{G})\)。我们有 \(\varphi*f\to f\) 在 \(\mathrm{L}^{p}(\mathrm{G})\) 中沿滤子 \(\mathfrak{B}\) 成立（INT，VIII，§4，第 7 节，命题 20）。此外，对 \(B_{0}\) 中每个 K 与 \(\varphi\in X_{K}\)，存在 \(\psi\in X_{K}^{\prime}\) 使 \(\varphi=\psi*\psi\)，从而 \(\varphi*f=\psi*(\psi*f)\)。我们有 \(\psi\in L^{1}(G)\cap L^{2}(G)\) 且 \(\psi*f\in L^{1}(G)\cap L^{2}(G)\)，从而 \(\varphi*f\in A(G)\)。

推论 1.　空间 A(G) 在 L \(^{1}\) (G) 与 L \(^{2}\) (G) 中稠密。它在 Stell(G) 中也稠密，并且它在 Fourier 变换下的像 \(\widehat{\mathrm{A}}(G)\) 在 \(\mathcal{C}_{0}(\widehat{\mathrm{G}})\) 中稠密。

命题的断言 (iv) 给出第一个断言。由于 L \(^{1}\) (G) 在 Stell(G) 中稠密，第二个断言由此得出，而最后一个断言进而由 II，第 209 页，命题 5 得出。

推论 2.　对 \(f \in \mathrm{A}(\mathrm{G})\)，设 \(\Omega_f\) 为由 \(\chi \in \widehat{\mathrm{G}}\)（满足 \(\mathcal{F}(f)(\chi) \neq 0\)）所成的集合。诸集合 \(\Omega_f\) 构成 \(\widehat{\mathrm{G}}\) 的一个开覆盖。

这由前一个推论得出，因为对每个 \(\chi \in \widehat{\mathrm{G}}\)，映射 \(f \mapsto \mathcal{F}(f)(\chi)\) 是 \(\mathrm{L}^1(\mathrm{G})\) 的一个非零特征标。

回忆 Stell(G) 的左正则表示 \(\pmb{\gamma}\)（在 \(\mathrm{L}^2 (\mathrm{G})\) 上；参见 I，第 125 页，第 13 节）记作 \(\pmb{\gamma}(\varphi)f = \varphi *f\)，其中 \(\varphi \in \mathrm{Stell}(\mathrm{G})\)，\(f\in \mathrm{L}^2 (\mathrm{G})\)。

引理 3.　对每个 \(f \in \mathrm{A}(\mathrm{G})\)，存在唯一的有界测度 \(\mu_{f}\)（在 \(\widehat{\mathrm{G}}\) 上），使得

\[
(\varphi * f) (e) = \int_ {\widehat {G}} \mathcal {F} (\varphi) d \mu_ {f}\tag{17}
\]

对所有 \(\varphi\in\mathrm{Stell}(\mathrm{G})\) 成立。

此外，对所有 \(f\) 与 \(g\)（在 A(G) 中），我们有等式

\[
\mathcal {F} (f) \cdot \mu_ {g} = \mathcal {F} (g) \cdot \mu_ {f},\tag{18}
\]

它是以 \(\mathcal{F}(f)\) 为密度（关于 \(\mu_g\)）的测度与以 \(\mathcal{F}(g)\) 为密度（关于 \(\mu_f\)）的测度之间的等式。

设 \(f\) 与 \(g\) 为 \(\mathrm{L}^{1}(\mathrm{G}) \cap \mathrm{L}^{2}(\mathrm{G})\) 的元素。对每个 \(\varphi \in \mathrm{Stell}(\mathrm{G})\)，有 \(\varphi * f \in \mathrm{L}^{2}(\mathrm{G})\) 且 \(\| \varphi * f \|_{2} \leqslant \| \varphi \|_{*} \| f \|_{2}\)（I，第 126 页，公式 (8)）。此外，我们有 \(\varphi * (f * g) = (\varphi * f) * g\)（同前，公式 (9)）。后一函数属于 \(\mathcal{C}_{0}(\mathrm{G})\)，即 \(\mathcal{K}(\mathrm{G})\) 在 \(\mathcal{C}(\mathrm{G})\) 中的闭包（INT，VIII，§4，第 5 节，命题 15）。此外，我们有

\[
\left\| \varphi * (f * g) \right\| _ {\infty} \leqslant \left\| \varphi * f \right\| _ {2} \| g \| _ {2} \leqslant \left\| \varphi \right\| _ {*} \| f \| _ {2} \| g \| _ {2}.
\]

由于诸函数 \(f * g\)（\(f\) 与 \(g\) 属于 \(\mathrm{L}^1(\mathrm{G}) \cap \mathrm{L}^2(\mathrm{G})\)）生成 \(\mathrm{A}(\mathrm{G})\)，由此推出 \(\varphi * f \in \mathcal{C}_0(\mathrm{G})\) 对所有 \(f \in \mathrm{A}(\mathrm{G})\) 与 \(\varphi \in \mathrm{Stell}(\mathrm{G})\) 成立，且映射 \(\varphi \mapsto (\varphi * f)(e)\) 是 Stell(G) 上的一个连续线性形式。由于 \(\mathcal{F}\) 是从 Stell(G) 到 \(\mathcal{C}_{0}(\widehat{\mathrm{G}})\) 上的同构（II，第 209 页，命题 5），第一个断言由此得出。

现设 \(f\) 与 \(g\) 属于 \(\mathrm{A}(\mathrm{G})\)。对 \(\varphi \in \mathrm{L}^1 (\mathrm{G})\) 有

\[
\begin{array}{r l r} & & {\big (\mathcal {F} (f) \cdot \mu_ {g} \big) (\mathcal {F} (\varphi)) = \int_ {\widehat {\mathrm{G}}} \mathcal {F} (\varphi) \mathcal {F} (f) d \mu_ {g} = \int_ {\widehat {\mathrm{G}}} \mathcal {F} (\varphi * f) d \mu_ {g}} \\ & & {= ((\varphi * f) * g) (e).} \end{array}\tag{19}
\]

由于 \((\varphi*f)*g=(\varphi*g)*f\)，且 \(L^{1}(G)\) 在 Fourier 变换下的像在 \(\mathcal{C}_{0}(\widehat{\mathrm{G}})\) 中稠密（II，第 210 页，推论），由公式 (19) 推出公式 (18) 对所有 \(f\) 与 \(g\)（在 \(A(G)\) 中）成立。

引理 4.　存在唯一的测度 \(\nu\)（在 \(\widehat{\mathbf{G}}\) 上），使得

\[
\mu_ {f} = \mathcal {F} (f) \cdot \nu
\]

对所有 \(f \in \mathrm{A}(\mathrm{G})\) 成立。对 \(f \in \mathrm{A}(\mathrm{G})\)，有 \(\mathcal{F}(f) \in \mathrm{L}^{1}(\widehat{\mathrm{G}}, \nu) \cap \mathrm{L}^{2}(\widehat{\mathrm{G}}, \nu)\)。设 \(f \in \mathrm{A}(\mathrm{G})\)。用 \(\Omega_{f}\) 表示 \(\widehat{G}\) 中由 \(\chi \in \widehat{G}\)（满足 \(\mathcal{F}(f)(\chi) \neq 0\)）组成的开集。设 \(\varphi\) 为 \(\widehat{G} - \Omega_{f}\) 的特征函数。于是对每个 \(g \in \mathrm{A}(\mathrm{G})\) 有

\[
\int_ {\widehat {\mathrm{G}}} \mathcal {F} (g) d (\boldsymbol {\varphi} \cdot \boldsymbol {\mu} _ {f}) = \int_ {\widehat {\mathrm{G}}} \varphi \mathcal {F} (f) d \mu_ {g} = 0
\]

这由公式 (18) 得到。

由 II，第 212 页，推论 1，像 \(\widehat{\mathrm{A}}(\mathrm{G})\)（\(\mathrm{A}(\mathrm{G})\) 在 Fourier 变换下的）在 \(\mathcal{C}_0(\widehat{\mathrm{G}})\) 中稠密。于是由前面的公式推出 \(\varphi \cdot \mu_f = 0\)，从而 \(\mu_f\) 集中在 \(\Omega_f\) 上（INT，IV，§4，第 7 节，定义 4）。设 \(\nu_f\) 为 \(\Omega_f\) 上以 \(\mathcal{F}(f)^{-1}\) 为密度关于 \(\mu_f|\Omega_f\) 的测度。

诸集合 \(\Omega_f\)（\(f \in \mathrm{A}(\mathrm{G})\)）构成 \(\widehat{\mathrm{G}}\) 的一个开覆盖（推论 2）。对所有 \(f\) 与 \(g\)（在 \(\mathrm{A}(\mathrm{G})\) 中），公式 (18) 表明 \(\nu_f|(\Omega_f \cap \Omega_g) = \nu_g|(\Omega_f \cap \Omega_g)\)。因此，存在唯一的测度 \(\nu\)（在 \(\widehat{\mathrm{G}}\) 上），使得 \(\nu_f = \nu|\Omega_f\) 对所有 \(f \in \mathrm{A}(\mathrm{G})\) 成立（INT，III，§2，第 1 节，命题 1）。

若 \(f \in \mathrm{A}(\mathrm{G})\)，测度 \(\mu_f\) 与 \(\mathcal{F}(f) \cdot \nu\) 都集中在 \(\Omega_f\) 上，并且它们在 \(\Omega_f\) 上的限制都等于 \(\mathcal{F}(f) \cdot \nu_f\)；因此这两个测度相等。

由于 \(\mu_f\) 是有界测度，Fourier 变换 \(\mathcal{F}(f)\) 属于空间 \(\mathrm{L}^1 (\widehat{\mathrm{G}},\nu)\)。此外，由于 \(\mathcal{F}(f)\) 属于 \(\mathcal{C}_0(\widehat{\mathrm{G}})\)，我们也有 \(\mathcal{F}(f)\in \mathrm{L}^2 (\widehat{\mathrm{G}},\nu)\)。

于是公式 (17) 可写成，对 \(\varphi\in\mathrm{Stell}(\mathrm{G})\) 与 \(f\in\mathrm{A}(\mathrm{G})\)：

\[
(\varphi * f) (e) = \int_ {\widehat {G}} \mathcal {F} (\varphi) \mathcal {F} (f) d \nu .\tag{20}
\]

特别地，对诸 \(f\) 与 \(g\)（在 \(\mathrm{A}(\mathrm{G})\) 中），有

\[
\int_ {\widehat {\mathrm{G}}} \mathcal {F} (f) \mathcal {F} (g) d \nu = (f * g) (e) = \int_ {\mathrm{G}} f (x) g (x ^ {- 1}) d x.\tag{21}
\]

命题 9.　引理 4 所刻画的测度 \(\nu\) 是 \(\widehat{\mathbf{G}}\) 上的一个 Haar 测度。

设 \(\chi \in \widehat{\mathrm{G}}\)。对诸 \(f\) 与 \(g\)（在 \(\mathrm{A}(\mathrm{G})\) 中），把公式 (21) 应用于 \(\chi f\) 与 \(\chi g\)。右边不变，故

\[
\nu (\mathcal {F} (f) \mathcal {F} (g)) = \nu (\mathcal {F} (\chi f) \mathcal {F} (\chi g)).
\]

由 II，第 208 页，公式 (12) 与 II，第 208 页，公式 (13)，可得

\[
\nu (\mathcal {F} (f) \mathcal {F} (g)) = \nu (\varepsilon_ {\chi} * (\mathcal {F} (f) \mathcal {F} (g))).
\]

然后由 INT，VIII，§4，第 3 节，命题 7 推出

\[
\nu (\mathcal {F} (f) \mathcal {F} (g)) = (\varepsilon_ {\chi^ {- 1}} * \nu) (\mathcal {F} (f) \mathcal {F} (g)),
\]

即

\[
(\mathscr {F} (f) \cdot \nu) (\mathscr {F} (g)) = (\mathscr {F} (f) \cdot (\varepsilon_ {\chi^ {- 1}} * \nu)) (\mathscr {F} (g)).
\]

由于空间 \(\widehat{\mathrm{A}}(\mathrm{G})\) 在 \(\mathcal{C}_{0}(\widehat{\mathrm{G}})\) 中稠密（推论 1），由此得等式

\[
\mathscr {F} (f) \cdot \nu = \mathscr {F} (f) \cdot (\varepsilon_ {\chi^ {- 1}} * \nu).
\]

因此测度 \(\nu\) 与 \(\varepsilon_{\chi^{-1}} * \nu\) 在开集 \(\Omega_f\)（\(\mathcal{F}(f)\) 在其上非零）上一致。由推论 2 以及 INT，III，§2，第 1 节，命题 1 的推论，这两个测度相等。

这表明测度 \(\nu\) 与 \(\widehat{\mathrm{G}}\) 上的某个 Haar 测度成比例。设 \(f \in \mathrm{A}(\mathrm{G})\)。在公式 (21) 中取 \(g = \widetilde{f}\)。它于是写成

\[
\int_ {\widehat {\mathrm{G}}} | \mathcal {F} (f) | ^ {2} d \nu = \int_ {\mathrm{G}} | f | ^ {2} d x,\tag{22}
\]

这就证明测度 \(\nu\) 非零。因此测度 \(\nu\) 是 \(\widehat{G}\) 上的一个 Haar 测度。

定义 4.　命题 9 中给出的 Haar 测度 \(\nu\)（在 \(\widehat{G}\) 上）称为 G 上给定的 Haar 测度 dx 的对偶测度。

我们常用 \(d\chi\) 或 \(d\widehat{x}\) 表示 \(\widehat{G}\) 上与 Haar 测度 dx 对偶的 Haar 测度。

注.　设 \(a\) 为实数 \(>0\)。若把 \(dx\) 换成测度 \(a \cdot dx\)，则函数 \(f\) 与 \(g \in \mathrm{L}^{1}(\mathrm{G})\) 的卷积换成 \(a(f*g)\)。我们已看到（II，第 209 页，注），\(\mathcal{F}(f)\) 换成 \(a\mathcal{F}(f)\)。因此 \(\mu_{f}\) 不变，而 \(\nu\) 换成 \(a^{-1} \cdot \nu\)。特别地，测度 \(dx \otimes d\hat{x}\)（在 \(\mathrm{G} \times \hat{\mathrm{G}}\) 上）与 \(\mathrm{G}\) 上 Haar 测度的选取无关。

引理 5.　空间 A(\(\hat{G}\)) 在 L²(\(\hat{G}\)) 中稠密。

设 \(h\) 为 \(\mathrm{L}^2(\widehat{\mathrm{G}})\) 中与 \(\widehat{\mathrm{A}}(\mathrm{G})\) 正交的一个元素。对诸 \(f\) 与 \(g\)（在 \(\mathrm{A}(\mathrm{G})\) 中），有 \(\mathcal{F}(f) \cdot \mathcal{F}(g) = \mathcal{F}(f*g) \in \widehat{\mathrm{A}}(\mathrm{G})\)，从而 \(h \cdot \overline{\mathcal{F}(f)}\) 与 \(\mathcal{F}(g)\) 正交。于是，对每个 \(f \in \mathrm{A}(\mathrm{G})\)，函数 \(h \cdot \mathcal{F}(f)\) 与 \(\widehat{\mathrm{A}}(\mathrm{G})\) 正交。但 \(h \cdot \mathcal{F}(f) \in \mathrm{L}^1(\widehat{\mathrm{G}})\)，且 \(\widehat{\mathrm{A}}(\mathrm{G})\) 在 \(\mathcal{C}_0(\widehat{\mathrm{G}})\) 中稠密，故测度 \(h\mathcal{F}(f) \cdot \nu\) 为零，即 \(h\mathcal{F}(f)\) 是 \(\nu\)-局部可忽略的（INT，V，§5，第 3 节，命题 3 的推论 2）。特别地，\(h\) 是 \(\nu\)-局部可忽略的，在集合 \(\Omega_f\)（即诸特征标 \(\chi\)（\(\mathcal{F}(f)(\chi) \neq 0\)）的集合）上。由推论 2，\(h\) 是 \(\nu\)-局部可忽略的，从而为零，因为 \(h\) 属于 \(\mathrm{L}^2(\widehat{\mathrm{G}})\)。证明完毕。

定理 1（Plancherel）.　Fourier 变换在 L²(G) 的子空间 A(G) 上的限制唯一地扩张为从 L²(G) 到 L²(Ĝ) 上的一个等距 Φ。

此外，若 \(f \in \mathrm{L}^{1}(\mathrm{G}) \cap \mathrm{L}^{2}(\mathrm{G})\)，则其 Fourier 变换属于 \(\mathrm{L}^{2}(\widehat{\mathrm{G}})\)，且在 \(\mathrm{L}^{2}(\widehat{\mathrm{G}})\) 中与 \(\Phi(f)\) 一致。

由公式 (22)，\(\mathcal{F}\) 在 A(G) 上的限制是子空间 A(G)（L \(^2\) (G) 的）到子空间 \(\hat{\mathrm{A}}(\mathrm{G})\)（L \(^2\) (\(\hat{\mathrm{G}}\)) 的）上的一个等距。由于 A(G) 在 L \(^2\) (G) 中稠密（II，第 212 页，推论 1），Fourier 变换唯一地扩张为等距 \(\Phi\)，它把 L \(^2\) (G) 映到 L \(^2\) (\(\hat{\mathrm{G}}\)) 的某个闭子空间上。但由于其像包含 \(\hat{\mathrm{A}}(\mathrm{G})\)（它在 L \(^2\) (\(\hat{\mathrm{G}}\)) 中稠密；引理 5），映射 \(\Phi\) 是满的。

现设 \(f \in \mathrm{L}^{1}(\mathrm{G}) \cap \mathrm{L}^{2}(\mathrm{G})\)；我们来证明其 Fourier 变换属于 \(\mathrm{L}^{2}(\widehat{\mathrm{G}})\)。由 II，第 211 页，命题 8, (iv) 以及 \(\mathrm{A}(\mathrm{G})\) 是 \(\mathrm{L}^{1}(\mathrm{G})\) 的理想这一事实，存在 \(\mathrm{A}(\mathrm{G})\) 上的一个滤子基 B，它在 \(\mathrm{L}^{1}(\mathrm{G})\) 与 \(\mathrm{L}^{2}(\mathrm{G})\) 中都收敛于 f。于是我们有

\[
\Phi (f) = \lim _ {g, \mathfrak {B}} \Phi (g) = \lim _ {g, \mathfrak {B}} \mathcal {F} (g)
\]

在 \(\mathrm{L}^{2}(\widehat{\mathrm{G}})\) 中成立，且 \(\mathcal{F}(f)=\lim_{g,\mathfrak{B}}\mathcal{F}(g)\) 在 \(\mathcal{C}_{0}(\widehat{\mathrm{G}})\) 中成立。因此存在一个序列 \((g_{n})\)（在 \(\mathrm{A}(\mathrm{G})\) 中），使 \(\mathcal{F}(g_{n})\) 收敛于 \(\Phi(f)\) 在 \(\mathrm{L}^{2}(\widehat{\mathrm{G}})\) 中，且收敛于 \(\mathcal{F}(f)\) 在 \(\mathcal{C}_{0}(\widehat{\mathrm{G}})\) 中。由 INT，IV，§3，第 4 节，定理 3 及其推论 1，有 \(\mathcal{F}(f)=\Phi(f)\)，特别地 \(\mathcal{F}(f)\in\mathrm{L}^{2}(\widehat{\mathrm{G}})\)。定理证毕。

我们仍用 \(\mathcal{F}\) 表示定理 1 中定义的从 \(\mathrm{L}^2 (\mathrm{G})\) 到 \(\mathrm{L}^2 (\widehat{\mathrm{G}})\) 上的等距，并称之为 \(\mathrm{L}^2 (\mathrm{G})\) 上的 Fourier 变换。同样，Fourier 余变换也唯一地等距扩张到 \(\mathrm{L}^2 (\mathrm{G})\) 上，仍称为 Fourier 余变换，记作 \(\overline{\mathcal{F}}\)。

推论.　假设 G 是紧的且 dx 是 G 上的规范化 Haar 测度。则 G 的酉特征标族是 L²(G) 的一组标准正交基。

由于 G 是紧的，G 的特征标属于 \(L^{2}(G)\)。对诸 \(\chi\) 与 \(\xi\)（在 \(\widehat{G}\) 中），有

\[
\int_ {\mathrm{G}} \overline {{\langle \chi , x \rangle}} \langle \xi , x \rangle d x = \mathcal {F} _ {\mathrm{G}} (d x) (\chi \xi^ {- 1}),
\]

因此 G 的特征标族是正交规范的（II，第 210 页，命题 6）。它还是完全的，因为 \(\chi \in \widehat{\mathbf{G}}\) 与 \(f \in \mathrm{L}^2(\mathrm{G})\) 的内积等于 \(\mathcal{F}_{\mathrm{G}}(f)(\chi)\)，从而 \(f\) 与 \(\widehat{\mathbf{G}}\) 正交当且仅当 \(\mathcal{F}_{\mathrm{G}}(f)\) 在 \(\mathrm{L}^2(\widehat{\mathbf{G}})\) 中为零，当且仅当 \(f\) 为零（定理 1）。

注.　1) 有关 \(\mathrm{L}^{1}(\mathrm{G})\) 上 Fourier 变换的一些公式可以推广到 \(\mathrm{L}^{2}(\mathrm{G})\) 上的 Fourier 变换。特别地，对 \(f \in \mathrm{L}^{2}(\mathrm{G})\) 与 \(\chi \in \widehat{\mathbf{G}}\)，有

\[
\overline {{\mathcal {F}}} (f) = (\chi \mapsto \mathcal {F} (f) (\chi^ {- 1})) = \overline {{\mathcal {F} (\overline {{f}})}}
\]

\[
\mathcal {F} (\varepsilon_ {x} * f) = \eta (x ^ {- 1}) \mathcal {F} (f), \qquad \overline {{\mathcal {F}}} (\varepsilon_ {x} * f) = \eta (x) \mathcal {F} (f)
\]

\[
\mathcal {F} (\chi f) = \varepsilon_ {\chi} * \mathcal {F} (f), \quad \overline {{\mathcal {F}}} (\chi f) = \varepsilon_ {\chi} * \mathcal {F} (f).
\]

若 \(\sigma\) 是 G 的自同构，\(\Delta\) 是 \(\sigma\) 的模（INT，VII，§1，第 4 节，定义 4），则对 \(f \in \mathrm{L}^2(\mathrm{G})\) 有

\[
\mathcal {F} (f \circ \sigma) = \Delta^ {- 1} \mathcal {F} (f) \circ \widehat {\sigma} ^ {- 1}
\]

在 L \(^{2}\) (G) 中成立。

\begin{enumerate}
\def\labelenumi{\arabic{enumi})}
\setcounter{enumi}{1}
\tightlist
\item
  公式
\end{enumerate}

\[
\| \overline {{\mathcal {F}}} (f) \| ^ {2} = \| f \| ^ {2}\tag{23}
\]

（\(f\in \mathrm{L}^2 (\mathrm{G})\)），或

\[
\int_ {\mathrm{G}} f (x) \overline {{g (x)}} d x = \int_ {\widehat {\mathrm{G}}} \mathcal {F} (f) (\chi) \overline {{\mathcal {F} (g) (\chi)}} d \chi\tag{24}
\]

（\(f\) 与 \(g\) 属于 \(\mathrm{L}^2 (\mathrm{G})\)）称为"Plancherel 公式"。

命题 10.　设 \(n \geqslant 0\) 为整数，\(\mathrm{G}_{1}, \cdots, \mathrm{G}_{n}\) 为局部紧交换群。设 \(\mu_{j}\)（\(1 \leqslant j \leqslant n\)）为 \(\mathrm{G}_{j}\) 上的 Haar 测度。设 G 为诸群 \(\mathrm{G}_{j}\)（\(1 \leqslant j \leqslant n\)）的乘积群。设 \(\beta\) 为 II，第 206 页，命题 2 中给出的从 \(\widehat{\mathrm{G}}\) 到 \(\prod \widehat{\mathrm{G}}_{j}\) 上的同构。\(\widehat{\mathrm{G}}\) 上与 G 上的乘积 Haar 测度 \(\mu = \mu_{1} \otimes \cdots \otimes \mu_{n}\) 对偶的 Haar 测度等同于诸 Haar 测度 \(\widehat{\mu}_{j}\) 的乘积。

事实上，对每个族 \((f_{j})\)（由 \(\mathcal{L}^{2}(\mathrm{G}_{j})\) 的非零元素组成），G 上由 \((x_{j})\mapsto\Pi f_{j}(x_{j})\) 定义的函数 f 属于 \(\mathcal{L}^{2}(\mathrm{G})\)，且满足

\[
\int_ {\mathrm{G}} | f | ^ {2} d \mu = \prod_ {j} \int_ {\mathrm{G} _ {j}} | f _ {j} | ^ {2} d \mu_ {j} = \prod_ {j} \int_ {\widehat {\mathrm{G}} _ {j}} | \mathcal {F} _ {\mathrm{G} _ {j}} (f _ {j}) | ^ {2} d \widehat {\mu} _ {j},
\]

这由 Plancherel 公式得出，它证明诸 \(\widehat{\mu}_{j}\) 的乘积 Haar 测度等同于 \(\mu\) 的对偶 Haar 测度。

\subsection{Fourier 反演公式}

回忆 A(G) 的每个元素 \(f\) 都是一个唯一的连续函数的类（II，第 210 页，命题 7, b)）。对 \(x \in \mathrm{G}\)，我们将用 \(f(x)\) 表示这个函数在 \(x\) 处的值。

命题 11.　设 \(f \in \mathrm{A}(\mathrm{G})\)。则 \(\mathcal{F}(f) \in \mathrm{L}^1(\widehat{\mathrm{G}})\)，且对每个 \(x \in \mathrm{G}\) 有

\[
f (x) = \int_ {\widehat {G}} \langle \widehat {x}, x \rangle \mathcal {F} (f) (\widehat {x}) d \widehat {x}.\tag{25}
\]

换句话说，对 \(f \in \mathrm{A}(\mathrm{G})\) 有

\[
f = \overline {{\mathcal {F}}} _ {\widehat {\mathrm{G}}} (\mathcal {F} _ {\mathrm{G}} (f)) \circ \eta ,\tag{26}
\]

其中 \(\eta\) 表示从 G 到二重对偶群 \(\widehat{\widehat{G}}\) 中的典范映射。

由 II，第 213 页，引理 4 与 II，第 214 页，命题 9，\(\mathcal{F}(f) \in \mathrm{L}^1(\widehat{\mathrm{G}})\) 对每个函数 \(f \in \mathrm{A}(\mathrm{G})\) 成立。由 Plancherel 公式 (24)，对诸 \(f\) 与 \(g\)（在 \(\mathrm{L}^2(\mathrm{G})\) 中）有

\[
(f * \widetilde {g}) (e) = \int_ {\widehat {\mathrm{G}}} \mathcal {F} (f) (\chi) \overline {{\mathcal {F} (g) (\chi)}} d \widehat {x} (\chi).\tag{27}
\]

设 \(f\) 与 \(g\) 属于 \(\mathrm{L}^1(\mathrm{G}) \cap \mathrm{L}^2(\mathrm{G})\)，并设 \(h = f * \widetilde{g} \in \mathrm{A}(\mathrm{G})\)。由于 Fourier 变换是一个对合态射，公式 (27) 就是断言 (25) 对函数 \(h\) 在点 \(x = e\) 处的情形。由线性性，我们推出公式 (25) 在点 \(x = e\) 处对每个函数 \(h \in \mathrm{A}(\mathrm{G})\) 成立。

设 \(x \in G\) 与 \(h \in A(G)\)。设 \(h_{1} = \varepsilon_{x^{-1}} * h\)。则 \(h_{1} \in A(G)\) 且 \(h_{1}(e) = h(x)\)。此外，由于 \(\mathcal{F}(h_{1})(\chi) = \overline{\langle\chi, x\rangle}\mathcal{F}(f)(\chi)\) 对所有 \(\chi \in \widehat{G}\) 成立（参见 II，第 208 页，公式 (11)），把公式 (25) 应用于函数 \(h_{1}\) 在点 e 处便得到公式 (25) 对 h 在点 x 处成立。

引理 6.　设 \(\varphi\in\mathrm{L}^{1}(\widehat{\mathrm{G}})\cap\mathrm{L}^{2}(\widehat{\mathrm{G}})\)。则 \(f=\overline{\mathcal{F}}_{\widehat{\mathrm{G}}}(\varphi)\circ\eta\) 属于 \(\mathrm{L}^{2}(\mathrm{G})\)，且 \(\mathcal{F}_{\mathrm{G}}(f)=\varphi\) 在 \(\mathrm{L}^{2}(\widehat{\mathrm{G}})\) 中成立。

函数 \(f\) 在 \(G\) 上连续且有界，因为 \(\varphi \in L^{1}(\widehat{G})\)。对每个函数 \(g \in L^{1}(G) \cap L^{2}(G)\) 有

\[
\begin{array}{r} \int_ {\mathrm{G}} g (x) f (x) d x = \int_ {\mathrm{G}} g (x) \Big (\int_ {\widehat {\mathrm{G}}} \langle \chi , x \rangle \varphi (\chi) d \widehat {x} (\chi) \Big) d x \\ = \int_ {\widehat {\mathrm{G}}} \overline {{\mathcal {F}}} _ {\mathrm{G}} (g) (\chi) \varphi (\chi) d \widehat {x} (\chi), \end{array}\tag{28}
\]

这是对函数 \((x,\chi)\mapsto g(x)\varphi(\chi)\langle\chi,x\rangle\) 应用 Lebesgue–Fubini 定理（INT，V，§8，第 4 节，定理 1, a)）得到的，该函数在 \(G\times\widehat{G}\) 上关于乘积测度 \(dx\otimes d\widehat{x}\) 可积。由此推出

\[
\left| \int_ {\mathrm{G}} g (x) f (x) d x \right| \leqslant \| \overline {{{\mathcal {F}}}} _ {\mathrm{G}} (g) \| _ {2} \| \varphi \| _ {2} = \| g \| _ {2} \| \varphi \| _ {2},
\]

这由 Plancherel 公式得出。于是线性形式 \(g \mapsto \int_{G} fg\) 在 \(\mathrm{L}^{1}(\mathrm{G}) \cap \mathrm{L}^{2}(\mathrm{G})\) 上连续，且由于 \(\mathrm{L}^{1}(\mathrm{G}) \cap \mathrm{L}^{2}(\mathrm{G})\) 在 Hilbert 空间 \(\mathrm{L}^{2}(\mathrm{G})\) 中稠密，我们推出 f 属于 \(\mathrm{L}^{2}(\mathrm{G})\)。

另一方面，应用 II，第 215 页，定理 1，我们得到

\[
\begin{array}{r} \int_ {\mathrm{G}} g (x) f (x) d x = \int_ {\widehat {\mathrm{G}}} \overline {{\mathcal {F} _ {\mathrm{G}} (\overline {{g}}) (\chi)}} \mathcal {F} _ {\mathrm{G}} (f) (\chi) d \widehat {x} (\chi) \\ = \int_ {\widehat {\mathrm{G}}} \overline {{\mathcal {F}}} _ {\mathrm{G}} (g) (\chi) \mathcal {F} _ {\mathrm{G}} (f) (\chi) d \widehat {x} (\chi) \end{array}
\]

对所有 \(g \in \mathrm{L}^{2}(\mathrm{G})\) 成立。与 (28) 比较，我们得出 \(\varphi = \mathcal{F}_{\mathrm{G}}(f)\) 在 \(\mathrm{L}^{2}(\widehat{\mathrm{G}})\) 中成立，因为 \(\mathrm{A}(\mathrm{G})\) 含于 \(\mathrm{L}^{1}(\mathrm{G}) \cap \mathrm{L}^{2}(\mathrm{G})\) 且 \(\widehat{\mathrm{A}}(\mathrm{G})\) 在 \(\mathrm{L}^{2}(\widehat{\mathrm{G}})\) 中稠密（II，第 215 页，引理 5）。

命题 12（Fourier 反演公式）

设 \(f \in \mathrm{L}^2(\mathrm{G})\) 满足 \(\mathcal{F}_{\mathrm{G}}(f) \in \mathrm{L}^1(\widehat{\mathrm{G}})\)。则 \(f = \overline{\mathcal{F}}_{\widehat{\mathrm{G}}}(\mathcal{F}_{\mathrm{G}}(f)) \circ \eta\) 在 \(\mathrm{L}^2(\mathrm{G})\) 中成立。换句话说，对几乎所有 \(x \in \mathrm{G}\) 有

\[
f (x) = \int_ {\widehat {G}} \langle \widehat {x}, x \rangle \mathcal {F} _ {G} (f) (\widehat {x}) d \widehat {x}.
\]

函数 \(\varphi = \mathcal{F}_{\mathrm{G}}(f)\) 属于 \(\mathrm{L}^1 (\widehat{\mathrm{G}})\cap \mathrm{L}^2 (\widehat{\mathrm{G}})\)，应用该引理即得所需公式。

推论 1.　对 \(\widehat{\mathbf{G}}\) 的每个闭子集 P 与每个不属于 P 的 \(\chi \in \widehat{\mathbf{G}}\)，存在一个函数 \(f \in \mathrm{L}^{1}(\mathrm{G})\)，使得 \(\mathcal{F}(f)\) 在 P 上为零而在 \(\chi\) 处非零。

由于由 (12) 有 \(\mathcal{F}(\chi f) = \varepsilon_{\chi} * \mathcal{F}(f)\) 对所有 \(\chi \in \widehat{\mathrm{G}}\)，只需考虑 \(\chi\) 是 \(\widehat{\mathrm{G}}\) 的单位元的情形。

设 U 为 \(e \in \widehat{G}\) 的一个紧对称邻域，使得 \(U^{2} \cap P = \varnothing\)。设 \(\varphi\) 为 \(\widehat{G}\) 上的一个连续正函数，在 U 外为零且满足 \(\varphi(e) = 1\)。于是函数 \(\varphi_{1} = \varphi * \varphi\) 在 P 上为零且 \(\varphi_{1}(e) > 0\)。因此只需证明 \(\varphi_{1}\) 属于 \(L^{1}(G)\) 上 Fourier 变换的像。此外，\(\varphi\) 与 \(\varphi_{1}\) 都属于 \(L^{1}(\widehat{G}) \cap L^{2}(\widehat{G})\)。令 \(f = \overline{\mathcal{F}}(\varphi) \circ \eta\) 与 \(f_{1} = \overline{\mathcal{F}}(\varphi_{1}) \circ \eta\)。引理 6 蕴含 f 与 \(f_{1}\) 属于 \(L^{2}(G)\) 且满足 \(\varphi = \mathcal{F}(f)\) 与 \(\varphi_{1} = \mathcal{F}(f_{1})\)。此外

\[
f _ {1} = \overline {{\mathcal {F}}} (\varphi * \varphi) \circ \eta = (\overline {{\mathcal {F}}} (\varphi) \circ \eta) ^ {2} = f ^ {2},
\]

从而 \(f_{1} \in \mathrm{L}^{1}(\mathrm{G})\)。于是 \(\varphi_{1} = \mathcal{F}(f_{1})\) 确实属于 \(\mathrm{L}^{1}(\mathrm{G})\) 在 Fourier 变换下的像。

推论 2.　Banach 代数 L \(^{1}\) (G) 是正则的（I，第 89 页，定义 1）。

由 I，第 88 页，命题 1 以及 L \(^{1}\) (G) 的 Gelfand 变换与 G 的 Fourier 余变换的等同，这由前一个推论得出。
